\section{The exact S6 cusp-to-gauge spectral chain}
\label{sec:s6-cusp-gauge-spectral-chain}

This section works on the regular complex two-torus fibres defined by the
displayed S6 period functions.  It does not use the source's later
identification of its compact total space with \(S^6\).  The metric, gauge
bundle, lattice discretization, and quantum oscillator introduced below are
all displayed explicitly; none is inferred from the phrase ``complex
torus.''  The source input used here is Definition~3.3, Theorem~3.4,
Lemma~3.14, and Proposition~3.21 of the exact transcription
\cite{S6Internal}.

\subsection{The radial cusp and the unchanged determinant}

Choose the source's distinguished cusp component and its positive radial
path
\begin{equation}\label{eq:s6cusp-radial-parameter}
 q_R:=e^{-2\pi R},\qquad R>R_0.
\end{equation}
Proposition~3.21 writes the period block as
\(Z=sB_0+C(t_c)\), with \(e^{2\pi i s}=t_c\), and proves that the three
functions \(h(t_c),\mu(t_c)\), and \(b(t_c)=\beta(t_c)+\tau(t_c)\) are
holomorphic at \(t_c=0\).  On the logarithm lift selected by
\eqref{eq:s6cusp-radial-parameter}, retain the source functions in the form
\begin{equation}\label{eq:s6cusp-source-functions}
 \tau_R=iR+h(q_R),\qquad
 \mu_R=\mu(q_R),\qquad
 \beta_R=b(q_R)-iR-h(q_R).
\end{equation}
Write all real and imaginary coordinates as
\begin{align}
 u_R&:=\operatorname{Re}\tau_R,
 &T_R&:=\operatorname{Im}\tau_R,
 &a_R&:=\operatorname{Re}\mu_R,
 &m_R&:=\operatorname{Im}\mu_R,\label{eq:s6cusp-real-imag-one}\\
 c_R^{\mathrm{re}}&:=\operatorname{Re}\beta_R,
 &p_R&:=\operatorname{Im}\beta_R,
 &h_R&:=\operatorname{Im}h(q_R),
 &d_R&:=\operatorname{Im}b(q_R).
 \label{eq:s6cusp-real-imag-two}
\end{align}
The identities in \eqref{eq:s6cusp-source-functions} give, without an
asymptotic replacement,
\begin{equation}\label{eq:s6cusp-imaginary-identities}
 T_R=R+h_R,\qquad m_R=\operatorname{Im}\mu(q_R),\qquad
 p_R=d_R-R-h_R=d_R-T_R.
\end{equation}

The real period matrix uses the ordered covering coordinates
\[
 (x^1,x^2,x^3,x^4)
 =
 (\operatorname{Re}\zeta_1,\operatorname{Im}\zeta_1,
  \operatorname{Re}\zeta_2,\operatorname{Im}\zeta_2).
\]
In that retained order it is
\begin{equation}\label{eq:s6cusp-real-period-matrix}
 P_R:=
 \begin{pmatrix}
  6a_R&u_R&1&0\\
  6m_R&T_R&0&0\\
  c_R^{\mathrm{re}}&a_R&0&1\\
  p_R&m_R&0&0
 \end{pmatrix}.
\end{equation}
Retain the positive cusp quantity
\begin{equation}\label{eq:s6cusp-delta}
 \delta_R:=6m_R^2-p_RT_R.
\end{equation}
The source's selected additive constant in \(\beta\) gives
\(T_R>0\) and
\(p_R-6m_R^2/T_R<0\), hence \(\delta_R>0\).

\begin{proposition}[Exact determinant and radial expansion]
\label{prop:s6cusp-determinant-expansion}
For every point of the radial cusp path,
\begin{align}
 \det P_R
 &=p_RT_R-6m_R^2=-\delta_R,
 \label{eq:s6cusp-det-period}\\
 \operatorname{Im}\tau_R
 \left(\operatorname{Im}\beta_R-
 \frac{6(\operatorname{Im}\mu_R)^2}
 {\operatorname{Im}\tau_R}\right)
 &=-\delta_R\ne0,
 \label{eq:s6cusp-unchanged-expression-delta}\\
 \delta_R
 &=T_R^2-d_RT_R+6m_R^2.
 \label{eq:s6cusp-delta-exact-radial}
\end{align}
Put
\begin{equation}\label{eq:s6cusp-boundary-values}
 h_0:=\operatorname{Im}h(0),\qquad
 m_0:=\operatorname{Im}\mu(0),\qquad
 d_0:=\operatorname{Im}b(0).
\end{equation}
Then
\begin{equation}\label{eq:s6cusp-delta-two-term-expansion}
 \delta_R=R^2+(2h_0-d_0)R+
 \bigl(h_0^2-d_0h_0+6m_0^2\bigr)
 +O\!\left(Re^{-2\pi R}\right).
\end{equation}
\end{proposition}

\begin{proof}
Expanding \eqref{eq:s6cusp-real-period-matrix} along its third and fourth
columns leaves the imaginary two-by-two block, with the sign retained:
\[
 \det P_R
 =\det\begin{pmatrix}6m_R&T_R\\p_R&m_R\end{pmatrix}
 =6m_R^2-p_RT_R
\]
multiplied by the sign of the column-row permutation used in the source's
real-coordinate order.  Direct expansion in that displayed order gives
\(p_RT_R-6m_R^2=-\delta_R\), agreeing with Lemma~3.14 of
\cite{S6Internal}.  Multiplication of the literal parenthesis in
\eqref{eq:s6cusp-unchanged-expression-delta} by
\(\operatorname{Im}\tau_R=T_R\) gives the same
 \(p_RT_R-6m_R^2\); division is legitimate because \(T_R>0\).

The last equality in \eqref{eq:s6cusp-imaginary-identities} gives
\[
 6m_R^2-p_RT_R
 =6m_R^2-(d_R-T_R)T_R
 =T_R^2-d_RT_R+6m_R^2,
\]
which is \eqref{eq:s6cusp-delta-exact-radial}.  Holomorphy at zero gives
\[
 h_R=h_0+O(e^{-2\pi R}),\quad
 m_R=m_0+O(e^{-2\pi R}),\quad
 d_R=d_0+O(e^{-2\pi R}).
\]
Substituting these three retained expansions and \(T_R=R+h_R\) into
\eqref{eq:s6cusp-delta-exact-radial}, and collecting every constant term,
gives \eqref{eq:s6cusp-delta-two-term-expansion}.
\end{proof}

\subsection{Marked flat torus, division lattice, and holonomy}

Define the regular fibre with the Euclidean covering metric by
\begin{equation}\label{eq:s6cusp-flat-fibre}
 F_R:=\R^4/P_R\Z^4,\qquad g_R:=\sum_{j=1}^4(dx^j)^2.
\end{equation}
The marking torus carries coordinates \(y\in\R^4/\Z^4\).

\begin{proposition}[Exact marking and finite division cellulation]
\label{prop:s6cusp-marking-cellulation}
The maps
\begin{align}
 \Phi_R:\R^4/\Z^4&\longrightarrow F_R,
 &[y]&\longmapsto[P_Ry],
 \label{eq:s6cusp-marking-map}\\
 \Phi_R^{-1}:F_R&\longrightarrow\R^4/\Z^4,
 &[x]&\longmapsto[P_R^{-1}x]
 \label{eq:s6cusp-marking-inverse}
\end{align}
are mutually inverse diffeomorphisms.  The pullback metric and volume are
\begin{equation}\label{eq:s6cusp-marking-metric-volume}
 G_R:=\Phi_R^*g_R=P_R^{\mathsf T}P_R,\qquad
 \operatorname{Vol}(F_R,g_R)=\sqrt{\det G_R}=\delta_R.
\end{equation}
For every integer \(N\ge2\), the map
\begin{equation}\label{eq:s6cusp-division-map}
 \iota_{R,N}:(\Z/N\Z)^4\longrightarrow F_R[N],\qquad
 [n]\longmapsto\left[\frac1N P_Rn\right]
\end{equation}
is a bijection onto the (N)-torsion subgroup.  Its coordinate edges are
\begin{equation}\label{eq:s6cusp-division-edges}
 e_{R,N,j}(n;t):=
 \Phi_R\!\left(\frac{n+te_j}{N}\right),\qquad0\le t\le1.
\end{equation}
\end{proposition}

\begin{proof}
If \(y-y'\in\Z^4\), then \(P_Ry-P_Ry'\in P_R\Z^4\), so
\eqref{eq:s6cusp-marking-map} is well defined.  If
\(x-x'\in P_R\Z^4\), then \(P_R^{-1}x-P_R^{-1}x'\in\Z^4\), so
\eqref{eq:s6cusp-marking-inverse} is well defined.  Their displayed
compositions are the identity before passage to either quotient.  For a
tangent vector \(v\) in \(y\)-coordinates,
\(\lVert d\Phi_R(v)\rVert^2=v^{\mathsf T}P_R^{\mathsf T}P_Rv\), proving
the metric formula.  The volume is
\(|\det P_R|=\delta_R\) by
\eqref{eq:s6cusp-det-period}.

Every element of \(F_R[N]\) has a representative \(x\) satisfying
\(Nx=P_Rz\) for some \(z\in\Z^4\), hence equals
\([P_Rz/N]\).  If \(\iota_{R,N}([n])=0\), then
\(P_Rn/N\in P_R\Z^4\).  Invertibility of \(P_R\) gives
\(n/N\in\Z^4\), so \([n]=0\) in \((\Z/N\Z)^4\).  This proves
surjectivity and injectivity.  Equation
\eqref{eq:s6cusp-division-edges} is then the image of each standard cubical
edge under the proved diffeomorphism.
\end{proof}

For a smooth connection \(A\) on a principal compact-\(G\) bundle which is
trivialized along the displayed cells, define its edge holonomy by
\begin{equation}\label{eq:s6cusp-holonomy-map}
 \operatorname{Hol}_{R,N}(A)_j(n):=
 \mathcal P\exp\!\left(
 -\int_0^1A_{e_{R,N,j}(n;t)}
 \left(\frac{P_Re_j}{N}\right)\,dt\right).
\end{equation}
For a gauge transformation \(g\), parallel transport gives exactly
\begin{equation}\label{eq:s6cusp-holonomy-gauge-equivariance}
 \operatorname{Hol}_{R,N}(A^g)_j(n)
 =g(\iota_{R,N}(n))^{-1}\operatorname{Hol}_{R,N}(A)_j(n)
 g(\iota_{R,N}(n+e_j)).
\end{equation}
Thus the continuum-to-link map has the usual finite gauge action as its
literal codomain action; this is an equivariant map, not an analogy.

\subsection{The complete dual map and the scalar spectrum}

For \(k=(k_1,k_2,k_3,k_4)\in\Z^4\), retain the real-part corrections
\begin{equation}\label{eq:s6cusp-adjusted-dual-integers}
 r_R(k):=k_1-6a_Rk_3-c_R^{\mathrm{re}}k_4,\qquad
 s_R(k):=k_2-u_Rk_3-a_Rk_4.
\end{equation}

\begin{theorem}[Exact dual-lattice isomorphism]
\label{thm:s6cusp-dual-isomorphism}
The maps
\begin{equation}\label{eq:s6cusp-dual-maps}
 \Z^4\longrightarrow(P_R\Z^4)^*,\quad
 k\longmapsto P_R^{-\mathsf T}k,\qquad
 q\longmapsto P_R^{\mathsf T}q
\end{equation}
are inverse group isomorphisms.  In the full displayed coordinates,
\begin{equation}\label{eq:s6cusp-full-dual-formula}
 P_R^{-\mathsf T}k=
 \left(
 k_3,
 \frac{m_Rr_R(k)-p_Rs_R(k)}{\delta_R},
 k_4,
 \frac{-T_Rr_R(k)+6m_Rs_R(k)}{\delta_R}
 \right).
\end{equation}
For (k=(r,s,0,0)), this becomes
\begin{equation}\label{eq:s6cusp-soft-dual-block}
 P_R^{-\mathsf T}(r,s,0,0)=
 \left(0,\frac{m_Rr-p_Rs}{\delta_R},0,
 \frac{-T_Rr+6m_Rs}{\delta_R}\right).
\end{equation}
The disappearance of \(u_R,a_R,c_R^{\mathrm{re}}\) from
\eqref{eq:s6cusp-soft-dual-block} follows from setting \(k_3=k_4=0\);
those real parts remain in the full inverse
\eqref{eq:s6cusp-full-dual-formula}.
\end{theorem}

\begin{proof}
A covector \(q\) pairs integrally with \(P_R\Z^4\) exactly when
\(P_R^{\mathsf T}q\in\Z^4\), proving the two inverse maps in
\eqref{eq:s6cusp-dual-maps}.  In the system \(P_R^{\mathsf T}q=k\), the
third and fourth column equations are
\(q_1=k_3\) and \(q_3=k_4\).  Substitution into the first two equations
leaves
\[
 \begin{pmatrix}6m_R&p_R\\T_R&m_R\end{pmatrix}
 \begin{pmatrix}q_2\\q_4\end{pmatrix}
 =\begin{pmatrix}r_R(k)\\s_R(k)\end{pmatrix}.
\]
The determinant of the displayed coefficient matrix is
\(6m_R^2-p_RT_R=\delta_R\), and its inverse is
\[
 \frac1{\delta_R}
 \begin{pmatrix}m_R&-p_R\\-T_R&6m_R\end{pmatrix}.
\]
This proves \eqref{eq:s6cusp-full-dual-formula} and its specialization.
\end{proof}

The character indexed by \(k\) is
\begin{equation}\label{eq:s6cusp-character}
 \chi_{R,k}([x]):=
 \exp\!\left(2\pi i\left\langle P_R^{-\mathsf T}k,x\right\rangle\right).
\end{equation}
For \(n\in\Z^4\), translation by \(P_Rn\) multiplies this expression by
\(e^{2\pi i k^{\mathsf T}n}=1\), so it is well defined.  The positive
scalar Laplacian satisfies
\begin{equation}\label{eq:s6cusp-character-eigenvalue}
 \Delta_R\chi_{R,k}=
 4\pi^2\lVert P_R^{-\mathsf T}k\rVert^2\chi_{R,k}.
\end{equation}

\begin{theorem}[Sharp scalar cusp spectrum and first correction]
\label{thm:s6cusp-sharp-scalar-spectrum}
For all sufficiently large \(R\), the shortest nonzero dual vectors are
indexed among
\((\pm1,0,0,0)\) and \((0,\pm1,0,0)\), and
\begin{equation}\label{eq:s6cusp-first-eigenvalue-correction}
 \lambda_1(F_R,g_R)=\frac{4\pi^2}{R^2}
 \left[
 1+\frac{\min\{2(d_0-h_0),-2h_0\}}{R}
 +O(R^{-2})
 \right].
\end{equation}
The exact two coordinate values before expansion are
\begin{align}
 \lVert P_R^{-\mathsf T}e_1\rVert^2
 &=\frac{m_R^2+T_R^2}{\delta_R^2},
 \label{eq:s6cusp-e1-exact}\\
 \lVert P_R^{-\mathsf T}e_2\rVert^2
 &=\frac{p_R^2+36m_R^2}{\delta_R^2}.
 \label{eq:s6cusp-e2-exact}
\end{align}
\end{theorem}

\begin{proof}
On the first two integer coordinates define
\begin{equation}\label{eq:s6cusp-dual-two-matrix}
 A_R:=\frac1{\delta_R}
 \begin{pmatrix}m_R&-p_R\\-T_R&6m_R\end{pmatrix}.
\end{equation}
Equations \eqref{eq:s6cusp-imaginary-identities} and
\eqref{eq:s6cusp-delta-two-term-expansion} give the operator-norm limit
\begin{equation}\label{eq:s6cusp-scaled-dual-limit}
 RA_R=
 \begin{pmatrix}0&1\\-1&0\end{pmatrix}+O_{\mathrm{op}}(R^{-1}).
\end{equation}
Consequently a constant \(C\) exists such that
\[
 \frac{1-C/R}{R}\lVert(r,s)\rVert
 \le\lVert A_R(r,s)\rVert
 \le\frac{1+C/R}{R}\lVert(r,s)\rVert
\]
for all \((r,s)\in\R^2\) and all large \(R\).  Integer vectors of
Euclidean norm at least \(\sqrt2\) are therefore longer than the two
coordinate vectors for all sufficiently large \(R\).  If \(k_3\ne0\) or
\(k_4\ne0\), \eqref{eq:s6cusp-full-dual-formula} gives
\[
 \lVert P_R^{-\mathsf T}k\rVert^2\ge k_3^2+k_4^2\ge1,
\]
whereas both coordinate values tend to zero.  This proves the asserted
finite list of minimizers.

Equations \eqref{eq:s6cusp-e1-exact} and
\eqref{eq:s6cusp-e2-exact} follow by putting \(e_1,e_2\) into
\eqref{eq:s6cusp-soft-dual-block}.  Direct substitution of
\eqref{eq:s6cusp-imaginary-identities} and
\eqref{eq:s6cusp-delta-two-term-expansion} gives
\begin{align*}
 \frac{m_R^2+T_R^2}{\delta_R^2}
 &=\frac1{R^2}\left[1+\frac{2(d_0-h_0)}R+O(R^{-2})\right],\\
 \frac{p_R^2+36m_R^2}{\delta_R^2}
 &=\frac1{R^2}\left[1-\frac{2h_0}R+O(R^{-2})\right].
\end{align*}
Taking their minimum and multiplying by \(4\pi^2\) proves
\eqref{eq:s6cusp-first-eigenvalue-correction}.
\end{proof}

\begin{proposition}[Low-energy lattice of modes]
\label{prop:s6cusp-low-energy-count}
Fix \(0<E<4\pi^2\), and count nonzero complex Fourier characters with
eigenvalue at most \(E\).  Then
\begin{equation}\label{eq:s6cusp-low-energy-count}
 N_R(E)=\frac{E}{4\pi}\delta_R+O_E(R).
\end{equation}
\end{proposition}

\begin{proof}
For large \(R\), the inequality \(E<4\pi^2\) and the exact lower bound on
modes with \(k_3\ne0\) or \(k_4\ne0\) restrict the count to
\(k=(r,s,0,0)\).  By \eqref{eq:s6cusp-dual-two-matrix}, these are the
integer points, except the origin, in
\[
 \left\{v\in\R^2:\lVert A_Rv\rVert^2\le\frac{E}{4\pi^2}\right\}.
\]
Since \(\det A_R=1/\delta_R\), this ellipse has area
\(E\delta_R/(4\pi)\).  Equation
\eqref{eq:s6cusp-scaled-dual-limit} makes both semiaxes \(O_E(R)\), so its
perimeter is \(O_E(R)\).  Covering each lattice point by its centered unit
square bounds the discrepancy from area by the area of a fixed-width
neighbourhood of the ellipse boundary, which is \(O_E(R)\).  Removing the
origin changes the count by one and proves the formula.
\end{proof}

\subsection{The mean-zero vertical direct integral and its continuous
threshold}

The radial parameter in \eqref{eq:s6cusp-radial-parameter} can itself be
retained as a noncompact measure variable.  This produces a precise
operator, rather than an inference from the collection of fibrewise
spectra.  Give every fibre its normalized Haar measure.  Under the marking
\(\Phi_R\), this measure is the Lebesgue probability measure on
\(\R^4/\Z^4\), so the characters
\[
 e_k(y):=e^{2\pi i k\cdot y},\qquad k\in\Z^4,
\]
provide one fixed measurable orthonormal basis along the entire radial
ray.  Put
\begin{equation}\label{eq:s6cusp-vertical-Hilbert-space}
 \mathscr H_{\mathrm v}:=
 L^2\!\left([R_0,\infty),\dd R;
 L^2_0(\R^4/\Z^4)\right),
\end{equation}
where the subscript zero deletes the character \(e_0\) on every fibre.

For \(k\ne0\), retain the complete fibre multiplier
\begin{equation}\label{eq:s6cusp-vertical-multiplier}
 \lambda_{R,k}:=4\pi^2
 \left\lVert P_R^{-\mathsf T}k\right\rVert^2.
\end{equation}
Define \(H_{\mathrm v}\) by its Fourier-coordinate domain and action:
\begin{align}
 \operatorname{Dom}(H_{\mathrm v})
 &:={}
 \left\{
 (f_k)_{k\in\Z^4\setminus\{0\}}:
 \sum_{k\ne0}\int_{R_0}^{\infty}
 \lambda_{R,k}^{2}|f_k(R)|^2\,\dd R<\infty
 \right\},
 \label{eq:s6cusp-vertical-domain}\\
 (H_{\mathrm v}f)_k(R)&:=\lambda_{R,k}f_k(R).
 \label{eq:s6cusp-vertical-action}
\end{align}
Equations \eqref{eq:s6cusp-vertical-domain}--
\eqref{eq:s6cusp-vertical-action} are the complete operator definition;
there is no derivative in the \(R\)-variable.

\begin{theorem}[Exact continuous threshold of the vertical operator]
\label{thm:s6cusp-vertical-continuous-threshold}
The operator \(H_{\mathrm v}\) is nonnegative and self-adjoint on
\(\mathscr H_{\mathrm v}\), and
\begin{equation}\label{eq:s6cusp-vertical-threshold}
 0\in\operatorname{Spec}_{\mathrm{ess}}(H_{\mathrm v})
 \cap\operatorname{Spec}_{\mathrm c}(H_{\mathrm v}),
 \qquad
 \ker H_{\mathrm v}=\{0\}.
\end{equation}
In particular, this explicitly defined mean-zero vertical operator has no
positive spectral gap.
\end{theorem}

\begin{proof}
The fibrewise Fourier transform and Fubini's theorem give the unitary map
and inverse
\begin{align*}
 \mathcal U:\mathscr H_{\mathrm v}
 &\longrightarrow
 \bigoplus_{k\in\Z^4\setminus\{0\}}
 L^2([R_0,\infty),\dd R),
 & (\mathcal Uf)_k(R)&:=\langle e_k,f(R)\rangle,\\
 \mathcal U^{-1}(f_k)(R,y)
 &:=\sum_{k\ne0}f_k(R)e_k(y).
\end{align*}
The inverse series converges in the displayed \(L^2\) space by Parseval.
Under \(\mathcal U\), equations \eqref{eq:s6cusp-vertical-domain}--
\eqref{eq:s6cusp-vertical-action} are the orthogonal sum of the maximal
multiplication operators by the real measurable functions
\(\lambda_{R,k}\).  A maximal real multiplication operator is
self-adjoint; the displayed square-integrability condition is exactly the
domain of their orthogonal sum.  Every multiplier is positive because
\(P_R\) is invertible and \(k\ne0\), proving nonnegativity.

For the fixed label \(k_*=(1,0,0,0)\), equation
\eqref{eq:s6cusp-e1-exact} and the radial estimates give constants
\(C_1,R_2>0\) such that
\begin{equation}\label{eq:s6cusp-vertical-e1-bound}
 0<\lambda_{R,k_*}\le \frac{C_1}{R^2}
 \qquad(R\ge R_2).
\end{equation}
Choose integers \(j\ge R_2\) and define
\begin{equation}\label{eq:s6cusp-vertical-Weyl-sequence}
 \Psi_j(R,y):=\mathbf1_{[j,j+1]}(R)e_{k_*}(y).
\end{equation}
Each \(\Psi_j\) has norm one, different \(j\)'s have disjoint supports,
and \eqref{eq:s6cusp-vertical-e1-bound} gives
\[
 \lVert H_{\mathrm v}\Psi_j\rVert^2
 =\int_j^{j+1}\lambda_{R,k_*}^{2}\,\dd R
 \le \frac{C_1^2}{j^4}\longrightarrow0.
\]
Thus \((\Psi_j)\) is an orthonormal Weyl sequence at zero.  Weyl's
criterion puts zero in the essential spectrum.

If \(H_{\mathrm v}f=0\), then
\(\lambda_{R,k}f_k(R)=0\) for almost every \(R\) and every \(k\ne0\).
The strict positivity of every \(\lambda_{R,k}\) forces every coefficient
\(f_k\) to vanish almost everywhere, so \(f=0\).  Hence zero is not an
eigenvalue.  A self-adjoint operator has empty residual spectrum, and zero
therefore lies in its continuous spectrum as asserted.
\end{proof}

The map from the direct integral to its fibres is evaluation in the
measurable-field sense, and the map \(\mathcal U\) above is its complete
Fourier diagonalization.  Since \(H_{\mathrm v}\) contains no horizontal
derivative and no boundary condition at the completed cusp, it is exactly a
vertical decomposable operator; equation
\eqref{eq:s6cusp-vertical-threshold} is asserted for that domain and no
other Hamiltonian.

\subsection{The flat Yang--Mills Hessian and its exact gauge quotient}

Let \(G\) be a compact Lie group, let \(\mathfrak g\) be its Lie algebra,
and fix an \(\operatorname{Ad}(G)\)-invariant positive inner product
\(\langle\ ,\ \rangle_{\mathfrak g}\).  On the trivial \(G\)-bundle over
\(F_R\), use
\begin{equation}\label{eq:s6cusp-yang-mills-action}
 \operatorname{YM}_R(A):=
 \frac{1}{2g_{\mathrm{YM}}^2}
 \int_{F_R}\langle F_A,F_A\rangle_{\mathfrak g}\,
 \operatorname{vol}_{g_R},
 \qquad g_{\mathrm{YM}}>0.
\end{equation}
Our curvature convention is
\[
 F_A=\dd A+A\wedge A,\qquad
 (A\wedge A)_{ij}=[A_i,A_j].
\]
Thus \(F_{ta}=t\,\dd a+t^2a\wedge a\), with every coefficient retained.

\begin{theorem}[Exact flat-connection quotient and Riesz Hessian]
\label{thm:s6cusp-flat-gauge-hessian}
At the flat connection \(A_0=0\), the infinitesimal gauge map, its kernel,
and its image are
\begin{equation}\label{eq:s6cusp-infinitesimal-gauge-map}
 \Gamma_R:C^\infty(F_R;\mathfrak g)
 \longrightarrow\Omega^1(F_R;\mathfrak g),
 \qquad
 \Gamma_R(\xi)=\dd\xi,
\end{equation}
\[
 \ker\Gamma_R=
 \{\text{constant \(\mathfrak g\)-valued functions}\},
 \qquad
 \operatorname{im}\Gamma_R=
 \{\dd\xi:\xi\in C^\infty(F_R;\mathfrak g)\}.
\]
The bilinear Hessian of \eqref{eq:s6cusp-yang-mills-action} is
\begin{equation}\label{eq:s6cusp-flat-yang-mills-bilinear-hessian}
 \operatorname{Hess}_{0}\operatorname{YM}_R[a,b]
 =\frac1{g_{\mathrm{YM}}^2}
 \langle\dd a,\dd b\rangle_{L^2\Omega^2}.
\end{equation}

Fix \(k\in\Z^4\setminus\{0\}\).  Write a complexified one-form on its
Fourier component in marking coordinates as
\[
 a_{k,\alpha,H}:=
 \chi_{R,k}\sum_{j=1}^4\alpha_j\,\dd y^j\otimes H,
 \qquad
 \alpha\in\C^4,\quad H\in\mathfrak g_{\C}.
\]
Define
\begin{equation}\label{eq:s6cusp-coulomb-projection}
 \Pi_{R,k}\alpha:=
 \alpha-k\,
 \frac{k^{\mathsf T}G_R^{-1}\alpha}
 {k^{\mathsf T}G_R^{-1}k}.
\end{equation}
Then
\begin{align}
 \ker\Pi_{R,k}&=\C k,\label{eq:s6cusp-coulomb-kernel}\\
 \operatorname{im}\Pi_{R,k}
 &=\{\alpha\in\C^4:
 k^{\mathsf T}G_R^{-1}\alpha=0\},
 \label{eq:s6cusp-coulomb-image}
\end{align}
and the mutually inverse maps
\begin{equation}\label{eq:s6cusp-coulomb-quotient-isomorphism}
 \C^4/\C k
 \underset{\alpha\mapsto[\alpha]}{\overset{[\alpha]\mapsto
 \Pi_{R,k}\alpha}{\rightleftarrows}}
 \operatorname{im}\Pi_{R,k}
\end{equation}
identify the Fourier-mode gauge quotient with the Coulomb slice.
On that slice the normalized operator and the Riesz Hessian operator are,
respectively,
\begin{align}
 \mathcal L_{R,k}
 &=4\pi^2 k^{\mathsf T}G_R^{-1}k\,I,
 \label{eq:s6cusp-normalized-flat-hessian}\\
 \mathcal H_{R,k}^{\mathrm{YM}}
 &=\frac{4\pi^2}{g_{\mathrm{YM}}^2}
 k^{\mathsf T}G_R^{-1}k\,I.
 \label{eq:s6cusp-riesz-flat-hessian}
\end{align}
\end{theorem}

\begin{proof}
The derivative at \(t=0\) of the gauge action
\(A\mapsto u^{-1}Au+u^{-1}\dd u\), along
\(u_t=\exp(t\xi)\), is \(\dd\xi\).  A function on the connected torus
has zero differential exactly when it is constant, proving the kernel and
image statements in \eqref{eq:s6cusp-infinitesimal-gauge-map}.

Substitution of \(F_{ta}=t\,\dd a+t^2a\wedge a\) into the literal action
\eqref{eq:s6cusp-yang-mills-action} gives
\[
 \operatorname{YM}_R(ta)
 =\frac1{2g_{\mathrm{YM}}^2}
 \left(
  t^2\lVert\dd a\rVert_{L^2}^2
  +2t^3\langle\dd a,a\wedge a\rangle_{L^2}
  +t^4\lVert a\wedge a\rVert_{L^2}^2
 \right).
\]
Polarization of the exact second derivative proves
\eqref{eq:s6cusp-flat-yang-mills-bilinear-hessian}.

In \(y\)-coordinates the inverse metric is \(G_R^{-1}\), so direct
differentiation and contraction give
\begin{equation}\label{eq:s6cusp-mode-divergence}
 \dd^*a_{k,\alpha,H}
 =-2\pi i\,\chi_{R,k}
 \bigl(k^{\mathsf T}G_R^{-1}\alpha\bigr)H.
\end{equation}
The Fourier coefficient of \(\dd(\chi_{R,k}H)\) is \(2\pi i kH\);
therefore the gauge line is exactly \(\C k\).
The denominator in \eqref{eq:s6cusp-coulomb-projection} is positive because
\(G_R^{-1}\) is positive definite and \(k\ne0\).  Direct substitution gives
\[
 \Pi_{R,k}k=0,\qquad
 k^{\mathsf T}G_R^{-1}\Pi_{R,k}\alpha=0,\qquad
 \Pi_{R,k}\alpha=\alpha
 \quad\text{on \eqref{eq:s6cusp-coulomb-image}.}
\]
These three identities prove the kernel, image, and both compositions in
\eqref{eq:s6cusp-coulomb-quotient-isomorphism}.

For completeness, put \(\omega=2\pi k\).  A direct coordinate contraction
gives
\begin{align*}
 \lVert\dd a_{k,\alpha,H}\rVert^2
 &=
 \left[
 (\omega^{\mathsf T}G_R^{-1}\omega)
 (\overline\alpha^{\mathsf T}G_R^{-1}\alpha)
 -
 \left|\omega^{\mathsf T}G_R^{-1}\alpha\right|^2
 \right]
 \lVert\chi_{R,k}H\rVert^2.
\end{align*}
The last term vanishes precisely on the Coulomb slice.  The Riesz
representative of \eqref{eq:s6cusp-flat-yang-mills-bilinear-hessian} is
therefore \(g_{\mathrm{YM}}^{-2}\mathcal L_{R,k}\), proving
\eqref{eq:s6cusp-normalized-flat-hessian} and
\eqref{eq:s6cusp-riesz-flat-hessian}.
\end{proof}

\begin{corollary}[Explicit non-gauge polarizations at the soft modes]
\label{cor:s6cusp-explicit-physical-soft-modes}
For \(k=(r,s,0,0)\ne0\), put
\[
 q_{R,r,s}:=P_R^{-\mathsf T}(r,s,0,0)
 =(0,q_2,0,q_4),
\]
where \(q_2,q_4\) are the two displayed components in
\eqref{eq:s6cusp-soft-dual-block}.  The maps
\begin{align}
 \mathcal J^{(1)}_{R,r,s}:\mathfrak g_{\C}
 &\longrightarrow\Omega^1(F_R;\mathfrak g_{\C}),
 &H&\longmapsto\chi_{R,(r,s,0,0)}\,\dd x^1\otimes H,
 \label{eq:s6cusp-soft-polarization-one}\\
 \mathcal J^{(3)}_{R,r,s}:\mathfrak g_{\C}
 &\longrightarrow\Omega^1(F_R;\mathfrak g_{\C}),
 &H&\longmapsto\chi_{R,(r,s,0,0)}\,\dd x^3\otimes H
 \label{eq:s6cusp-soft-polarization-three}
\end{align}
are injective, lie in the Coulomb slice, have zero intersection with the
gauge line, and are eigenmaps of \(\mathcal L_R\) with eigenvalue
\begin{equation}\label{eq:s6cusp-soft-gauge-eigenvalue}
 4\pi^2\lVert q_{R,r,s}\rVert^2.
\end{equation}
After the harmonic \(k=0\) one-forms are removed, the first positive
normalized gauge-Hessian eigenvalue is exactly
\(\lambda_1(F_R,g_R)\), and the first Riesz-Hessian eigenvalue is
\[
 \frac1{g_{\mathrm{YM}}^2}\lambda_1(F_R,g_R).
\]
Each minimizing unordered Fourier pair \(\{k,-k\}\) contributes
\(6\dim\mathfrak g\) real dimensions: two trigonometric modes, three
transverse polarizations, and \(\dim\mathfrak g\) Lie-algebra directions.
\end{corollary}

\begin{proof}
In covering \(x\)-coordinates,
\[
 \dd^*(\chi_{R,k}\dd x^jH)
 =-2\pi i(q_{R,k})_j\chi_{R,k}H.
\]
Equation \eqref{eq:s6cusp-soft-dual-block} gives
\((q_{R,r,s})_1=(q_{R,r,s})_3=0\), proving the two Coulomb identities.
The same equation and \(k\ne0\) imply \(q_{R,r,s}\ne0\).  Its gauge line
is \(\C q_{R,r,s}\subset\operatorname{span}_{\C}\{\dd x^2,\dd x^4\}\);
it consequently intersects each image in
\eqref{eq:s6cusp-soft-polarization-one}--
\eqref{eq:s6cusp-soft-polarization-three} only at zero.  Constant
coefficient one-forms on a flat covering have the scalar Fourier
eigenvalue, which proves \eqref{eq:s6cusp-soft-gauge-eigenvalue}.
For every nonzero \(k\), the Coulomb fibre has complex dimension three
over each Lie-algebra direction.  Pairing \(k\) with \(-k\) converts one
complex character coefficient into its two real sine/cosine coefficients,
which proves the stated real dimension and the first-eigenvalue formulas.
\end{proof}

\subsection{The exact metric cubical Hessian and its Gaussian quantization}

Fix \(N\ge2\) and write \(V_N=(\Z/N\Z)^4\).  A zero-cochain is a function
\(\xi:V_N\to\mathfrak g\), and a one-cochain is an array
\(a=(a_j(n))_{n\in V_N,1\le j\le4}\).  Define
\begin{align}
 (d_N^0\xi)_j(n)
 &:=N\bigl(\xi(n+e_j)-\xi(n)\bigr),
 \label{eq:s6cusp-discrete-d-zero}\\
 (d_N^1a)_{ij}(n)
 &:=N\bigl(
 a_j(n+e_i)-a_j(n)
 -a_i(n+e_j)+a_i(n)
 \bigr).
 \label{eq:s6cusp-discrete-d-one}
\end{align}
The metric cochain inner products are
\begin{align}
 \langle a,b\rangle_{1,R,N}
 &:=
 \frac{\sqrt{\det G_R}}{N^4}
 \sum_{n\in V_N}
 (G_R^{-1})^{ij}
 \langle a_i(n),b_j(n)\rangle_{\mathfrak g},
 \label{eq:s6cusp-one-cochain-inner-product}\\
 \langle F,K\rangle_{2,R,N}
 &:=
 \frac{\sqrt{\det G_R}}{2N^4}
 \sum_{n\in V_N}
 (G_R^{-1})^{ia}(G_R^{-1})^{jb}
 \langle F_{ij}(n),K_{ab}(n)\rangle_{\mathfrak g}.
 \label{eq:s6cusp-two-cochain-inner-product}
\end{align}
Both retain the full matrix \(G_R=P_R^{\mathsf T}P_R\).

For \(k\in V_N\), define
\begin{equation}\label{eq:s6cusp-discrete-fourier-symbol}
 \zeta_{N,k}(n):=e^{2\pi i k\cdot n/N},
 \qquad
 \delta_{N,j}(k):=N\left(e^{2\pi i k_j/N}-1\right).
\end{equation}
The complexified Fourier one-cochain
\(a_j(n)=\alpha_j\zeta_{N,k}(n)H\) then satisfies
\begin{equation}\label{eq:s6cusp-discrete-curvature-symbol}
 (d_N^1a)_{ij}(n)=
 \bigl(\delta_{N,i}(k)\alpha_j
 -\delta_{N,j}(k)\alpha_i\bigr)
 \zeta_{N,k}(n)H.
\end{equation}

\begin{theorem}[Exact discrete gauge quotient and cochain spectrum]
\label{thm:s6cusp-discrete-gauge-spectrum}
For \(k\ne0\) in \(V_N\), put \(\delta=\delta_N(k)\) and define
\begin{equation}\label{eq:s6cusp-discrete-coulomb-projection}
 \Pi_{R,N,k}\alpha:=
 \alpha-\delta\,
 \frac{\overline\delta^{\mathsf T}G_R^{-1}\alpha}
 {\overline\delta^{\mathsf T}G_R^{-1}\delta}.
\end{equation}
Its exact kernel, image, quotient map, and inverse are
\begin{align}
 \ker\Pi_{R,N,k}&=\C\delta,
 \label{eq:s6cusp-discrete-coulomb-kernel}\\
 \operatorname{im}\Pi_{R,N,k}
 &=\{\alpha:\overline\delta^{\mathsf T}G_R^{-1}\alpha=0\},
 \label{eq:s6cusp-discrete-coulomb-image}\\
 \C^4/\C\delta
 &\underset{\alpha\mapsto[\alpha]}{\overset{[\alpha]\mapsto
 \Pi_{R,N,k}\alpha}{\rightleftarrows}}
 \operatorname{im}\Pi_{R,N,k}.
 \label{eq:s6cusp-discrete-coulomb-quotient}
\end{align}
The Riesz operator \(L_{R,N}\) of the quadratic form
\(\langle d_N^1a,d_N^1a\rangle_{2,R,N}\), restricted to the nonzero
discrete Coulomb modes, is scalar on the \(k\)-th transverse fibre with
eigenvalue
\begin{equation}\label{eq:s6cusp-discrete-hessian-eigenvalue}
 \lambda_{R,N}(k)=
 \overline{\delta_N(k)}^{\mathsf T}
 G_R^{-1}\delta_N(k).
\end{equation}
\end{theorem}

\begin{proof}
The Fourier coefficient of \(d_N^0\xi\) is \(\delta\,\widehat\xi(k)\), so
the discrete gauge line is exactly \(\C\delta\).  Since at least one
component of \(k\) is nonzero modulo \(N\), one has \(\delta\ne0\); positive
definiteness of \(G_R^{-1}\) makes the denominator in
\eqref{eq:s6cusp-discrete-coulomb-projection} positive.  The three direct
substitutions used after \eqref{eq:s6cusp-coulomb-projection} apply with
\(k\) replaced by \(\delta\) and the transpose replaced by the Hermitian
transpose.  They prove \eqref{eq:s6cusp-discrete-coulomb-kernel}--
\eqref{eq:s6cusp-discrete-coulomb-quotient}.

For arbitrary \(\delta,\alpha\in\C^4\), expansion of every index in
\eqref{eq:s6cusp-two-cochain-inner-product} gives the exact identity
\begin{align}
 &\frac12
 (\overline\delta_i\overline\alpha_j
  -\overline\delta_j\overline\alpha_i)
 (G_R^{-1})^{ia}(G_R^{-1})^{jb}
 (\delta_a\alpha_b-\delta_b\alpha_a)
 \notag\\
 &\qquad=
 (\overline\delta^{\mathsf T}G_R^{-1}\delta)
 (\overline\alpha^{\mathsf T}G_R^{-1}\alpha)
 -
 \left|\overline\delta^{\mathsf T}G_R^{-1}\alpha\right|^2.
 \label{eq:s6cusp-discrete-wedge-identity}
\end{align}
The last term vanishes exactly on
\eqref{eq:s6cusp-discrete-coulomb-image}.  Comparison with
\eqref{eq:s6cusp-one-cochain-inner-product} proves that the Riesz operator
is multiplication by \eqref{eq:s6cusp-discrete-hessian-eigenvalue}.
\end{proof}

Put
\begin{equation}\label{eq:s6cusp-discrete-sN}
 s_N:=2N\sin\left(\frac{\pi}{N}\right).
\end{equation}

\begin{theorem}[Sharp finite-cubical cusp scale]
\label{thm:s6cusp-sharp-discrete-cusp-scale}
There is \(R_1\) independent of \(N\) such that, for all \(R\ge R_1\) and
all \(N\ge2\), the first positive cochain-Hessian eigenvalue is
\begin{equation}\label{eq:s6cusp-first-discrete-exact}
 \lambda_{1,R,N}^{\mathrm{coch}}
 =s_N^2\min\left\{
 \frac{m_R^2+T_R^2}{\delta_R^2},
 \frac{p_R^2+36m_R^2}{\delta_R^2}
 \right\}.
\end{equation}
The realizing residues are among
\(\pm e_1,\pm e_2\), with \(e_j=-e_j\) when \(N=2\).  Equivalently,
\begin{equation}\label{eq:s6cusp-first-discrete-expansion}
 \lambda_{1,R,N}^{\mathrm{coch}}
 =\frac{s_N^2}{R^2}
 \left[
 1+\frac{\min\{2(d_0-h_0),-2h_0\}}{R}
 +O(R^{-2})
 \right],
\end{equation}
where the \(O(R^{-2})\) constant is independent of \(N\).  For
\(j=1,2\), the coordinate-mode continuum and discrete eigenvalues satisfy
the exact identity
\begin{equation}\label{eq:s6cusp-continuum-discrete-ratio}
 \frac{\lambda_{R,N}(e_j)}
 {4\pi^2\lVert P_R^{-\mathsf T}e_j\rVert^2}
 =
 \left[
 \frac{N\sin(\pi/N)}{\pi}
 \right]^2.
\end{equation}
\end{theorem}

\begin{proof}
Since \(G_R^{-1}=P_R^{-1}P_R^{-\mathsf T}\),
\[
 \lambda_{R,N}(k)
 =\lVert P_R^{-\mathsf T}\delta_N(k)\rVert^2.
\]
For \(k=e_j\), the vector \(\delta_N(k)\) is
\(N(e^{2\pi i/N}-1)e_j\), whose squared modulus is \(s_N^2\).
Equations \eqref{eq:s6cusp-e1-exact} and
\eqref{eq:s6cusp-e2-exact} therefore give the two entries in
\eqref{eq:s6cusp-first-discrete-exact}, and comparison with the factor
\(4\pi^2\) in \eqref{eq:s6cusp-character-eigenvalue} proves
\eqref{eq:s6cusp-continuum-discrete-ratio}.

For a nonzero residue \(\ell\in\Z/N\Z\),
\[
 \left|N(e^{2\pi i\ell/N}-1)\right|\ge s_N.
\]
If \(\ell\not\equiv\pm1\pmod N\), the left side is at least
\(\sqrt2\,s_N\); for \(N=2,3\) there is no additional residue class.
If \(k_3\ne0\) or \(k_4\ne0\), the first and third components of
\eqref{eq:s6cusp-full-dual-formula}, with the formula extended
\(\C\)-linearly, give
\[
 \lambda_{R,N}(k)\ge
 |\delta_{N,3}(k)|^2+|\delta_{N,4}(k)|^2\ge s_N^2.
\]
For \(k_3=k_4=0\), equation
\eqref{eq:s6cusp-scaled-dual-limit} gives constants \(C,R_0'\),
independent of \(N\), such that
\[
 \frac{1-C/R}{R}\lVert(\delta_{N,1},\delta_{N,2})\rVert
 \le
 \lVert A_R(\delta_{N,1},\delta_{N,2})\rVert
 \le
 \frac{1+C/R}{R}\lVert(\delta_{N,1},\delta_{N,2})\rVert.
\]
Every soft residue other than the four stated coordinate residues has
Euclidean norm at least \(\sqrt2\,s_N\).  For one fixed sufficiently large
\(R_1\), the lower bound for those residues exceeds the upper bound for
the coordinate residues, while \(s_N^2\) exceeds every coordinate value.
This proves the uniform minimizing list.  The two expansions in the proof
of Theorem~\ref{thm:s6cusp-sharp-scalar-spectrum}, multiplied by the bounded
factor \(s_N^2\), prove \eqref{eq:s6cusp-first-discrete-expansion} with a
constant independent of \(N\).
\end{proof}

\begin{proposition}[Exact physical mesh criterion]
\label{prop:s6cusp-physical-mesh-criterion}
The largest Euclidean edge length of the division cellulation is
\begin{equation}\label{eq:s6cusp-physical-mesh-exact}
 \operatorname{mesh}(R,N)=
 \frac1N\max_{1\le j\le4}\lVert P_Re_j\rVert.
\end{equation}
Along any sequence \(R\to\infty\), \(N=N(R)\ge2\),
\begin{equation}\label{eq:s6cusp-physical-mesh-asymptotic}
 \operatorname{mesh}(R,N(R))
 =\frac{R}{N(R)}\left[1+O(R^{-1})\right].
\end{equation}
Consequently the physical cells shrink to zero exactly when
\(N(R)/R\to\infty\).  This geometric criterion is distinct from the
coordinate-mode spectral ratio \eqref{eq:s6cusp-continuum-discrete-ratio},
which tends to one whenever \(N\to\infty\).
\end{proposition}

\begin{proof}
Equation \eqref{eq:s6cusp-division-edges} has constant tangent vector
\(P_Re_j/N\), proving \eqref{eq:s6cusp-physical-mesh-exact}.  The four
column lengths are
\begin{align*}
 \lVert P_Re_1\rVert^2
 &=(6a_R)^2+(6m_R)^2+(c_R^{\mathrm{re}})^2+p_R^2,\\
 \lVert P_Re_2\rVert^2
 &=u_R^2+T_R^2+a_R^2+m_R^2,\\
 \lVert P_Re_3\rVert^2&=1,\qquad
 \lVert P_Re_4\rVert^2=1.
\end{align*}
Holomorphy at \(q_R=0\), together with
\(p_R=-R+(d_0-h_0)+O(e^{-2\pi R})\) and
\(T_R=R+h_0+O(e^{-2\pi R})\), shows that each of the first two lengths is
\(R[1+O(R^{-1})]\).  This proves
\eqref{eq:s6cusp-physical-mesh-asymptotic} and both implications of the
stated criterion.  The final spectral statement follows directly from
\(N\sin(\pi/N)/\pi\to1\).
\end{proof}

Let \(\mathscr C_{R,N}\) denote the real, nonharmonic discrete Coulomb
space: it is the orthogonal direct sum, with the reality condition coupling
\(k\) to \(-k\), of the fibres
\eqref{eq:s6cusp-discrete-coulomb-image} over \(k\ne0\).
The operator \(L_{R,N}\) is strictly positive on this finite-dimensional
space.

\begin{theorem}[Exact Gaussian quantization of the cochain Hessian]
\label{thm:s6cusp-gaussian-quantization}
Equip \(\mathscr C_{R,N}\) with the Lebesgue measure determined by
\eqref{eq:s6cusp-one-cochain-inner-product}.  The explicitly defined
quadratic Hamiltonian
\begin{equation}\label{eq:s6cusp-quadratic-hamiltonian}
 \widehat H^{\mathrm{quad}}_{R,N}
 :=
 -\frac{g_{\mathrm{YM}}^2}{2}\Delta_{\mathscr C_{R,N}}
 +\frac1{2g_{\mathrm{YM}}^2}
 \langle a,L_{R,N}a\rangle_{1,R,N}
\end{equation}
is essentially self-adjoint on the Schwartz space and has excitation gap
\begin{equation}\label{eq:s6cusp-quadratic-gap-exact}
 \operatorname{gap}
 \bigl(\widehat H^{\mathrm{quad}}_{R,N}\bigr)
 =\sqrt{\lambda_{1,R,N}^{\mathrm{coch}}}.
\end{equation}
For all \(R\ge R_1\),
\begin{equation}\label{eq:s6cusp-quadratic-gap-asymptotic}
 \operatorname{gap}
 \bigl(\widehat H^{\mathrm{quad}}_{R,N}\bigr)
 =
 \frac{s_N}{R}
 \left[
 1+\frac{\min\{d_0-h_0,-h_0\}}{R}
 +O(R^{-2})
 \right].
\end{equation}
In particular the gap tends to zero for fixed \(N\), and also for every
sequence \(N=N(R)\ge2\), because
\begin{equation}\label{eq:s6cusp-sN-uniform-bounds}
 4\le s_N<2\pi.
\end{equation}
\end{theorem}

\begin{proof}
Choose an orthonormal real eigenbasis of the positive symmetric operator
\(L_{R,N}\), with eigenvalues \(\lambda_\nu>0\).  This is an orthogonal
coordinate isomorphism with inverse given by expansion in the same basis;
under it, \eqref{eq:s6cusp-quadratic-hamiltonian} is the sum of the
one-dimensional operators
\[
 H_{\lambda_\nu,g_{\mathrm{YM}}}
 =-\frac{g_{\mathrm{YM}}^2}{2}\frac{\dd^2}{\dd x_\nu^2}
 +\frac{\lambda_\nu}{2g_{\mathrm{YM}}^2}x_\nu^2.
\]
For each \(\lambda>0\), the unitary map and its inverse are
\begin{align*}
 (U_{\lambda,g_{\mathrm{YM}}}f)(u)
 &:=
 \left(\frac{g_{\mathrm{YM}}}{\lambda^{1/4}}\right)^{1/2}
 f\left(\frac{g_{\mathrm{YM}}}{\lambda^{1/4}}u\right),\\
 (U_{\lambda,g_{\mathrm{YM}}}^{-1}\varphi)(x)
 &:=
 \left(\frac{\lambda^{1/4}}{g_{\mathrm{YM}}}\right)^{1/2}
 \varphi\left(\frac{\lambda^{1/4}}{g_{\mathrm{YM}}}x\right).
\end{align*}
Direct differentiation, with no coefficient suppressed, gives
\[
 U_{\lambda,g_{\mathrm{YM}}}
 H_{\lambda,g_{\mathrm{YM}}}
 U_{\lambda,g_{\mathrm{YM}}}^{-1}
 =
 \sqrt\lambda\left(
 -\frac12\frac{\dd^2}{\dd u^2}+\frac12u^2
 \right).
\]
The Hermite basis therefore gives energy increments
\(\sqrt{\lambda_\nu}\); the smallest increment proves
\eqref{eq:s6cusp-quadratic-gap-exact}.  Taking the positive square root of
\eqref{eq:s6cusp-first-discrete-expansion} gives
\eqref{eq:s6cusp-quadratic-gap-asymptotic}.  Finally,
\(\sin x\ge2x/\pi\) for \(0\le x\le\pi/2\) and
\(\sin x<x\) for \(x>0\) give \eqref{eq:s6cusp-sN-uniform-bounds}.
\end{proof}

\subsection{The native degree-six class on the full four-torus}
\label{subsec:s6cusp-native-four-dimensional-magnetic}

The second-pass magnetic calculation becomes intrinsic to the marked S6
fibre once the source class is retained.  Let
\begin{equation}\label{eq:s6cusp-integral-coframe}
 \vartheta_R^j:=(P_R^{-\mathsf T}e_j)\mathbin{\cdot}\dd x,
 \qquad 1\le j\le4.
\end{equation}
Since \(x=P_Ry\), one has \(\vartheta_R^j=\dd y^j\), and therefore
\begin{equation}\label{eq:s6cusp-integral-periods}
 \int_{P_Re_k}\vartheta_R^j=\delta_k^j.
\end{equation}
In the source identification
\((\dd y^1,\dd y^2,\dd y^3,\dd y^4)=(\gamma,u,w,\delta)\),
the monodromy-invariant integral \((1,1)\)-class of Lemma~9.4 is
\begin{equation}\label{eq:s6cusp-native-eta}
 \eta_R:=
 \vartheta_R^2\wedge\vartheta_R^3
 +6\vartheta_R^1\wedge\vartheta_R^4
 =u\wedge w+6\gamma\wedge\delta
 \in H^2(F_R;\Z)
 \end{equation}
\cite{S6Internal}.  Its alternating matrix in the ordered source basis is
\begin{equation}\label{eq:s6cusp-native-eta-matrix}
 E_\eta:=
 \begin{pmatrix}
 0&0&0&6\\
 0&0&1&0\\
 0&-1&0&0\\
 -6&0&0&0
 \end{pmatrix}.
\end{equation}
Direct exterior multiplication retains the orientation and gives
\begin{equation}\label{eq:s6cusp-native-eta-square-pfaffian}
 \frac12\eta_R\wedge\eta_R
 =6\,\dd y^1\wedge\dd y^2\wedge\dd y^3\wedge\dd y^4,
 \qquad
 \operatorname{Pf}(E_\eta)=6.
\end{equation}

\begin{proposition}[Generator-by-generator monodromy invariance]
\label{prop:s6cusp-native-eta-monodromy}
Let the source monodromy generators act on
\(V=\Z\langle\gamma,u,w,\delta\rangle\) by
\begin{equation}\label{eq:s6cusp-source-cohomology-monodromy}
 T_1=
 \begin{pmatrix}
 1&0&-6&2\\
 0&-1&1&1\\
 0&-1&0&1\\
 0&0&0&1
 \end{pmatrix},
 \qquad
 T_2=
 \begin{pmatrix}
 1&6&0&-3\\
 0&0&-1&1\\
 0&1&0&0\\
 0&0&0&1
 \end{pmatrix},
\end{equation}
where columns are images in the ordered basis
\((\gamma,u,w,\delta)\).  Then
\begin{equation}\label{eq:s6cusp-native-eta-generator-invariance}
 (\bigwedge^2T_1)\eta_R=\eta_R,
 \qquad
 (\bigwedge^2T_2)\eta_R=\eta_R.
\end{equation}
Consequently the rank-one subgroup \(\Z\eta_R\) is carried identically
around every loop in the regular base.
\end{proposition}

\begin{proof}
Reading the columns of \eqref{eq:s6cusp-source-cohomology-monodromy} gives
\begin{align*}
 T_1\gamma&=\gamma,&
 T_1u&=-u-w,&
 T_1w&=-6\gamma+u,&
 T_1\delta&=2\gamma+u+w+\delta,\\
 T_2\gamma&=\gamma,&
 T_2u&=6\gamma+w,&
 T_2w&=-u,&
 T_2\delta&=-3\gamma+u+\delta.
\end{align*}
Exterior multiplication, without changing the source basis, yields
\begin{align*}
 (\bigwedge^2T_1)\eta_R
 &=(-u-w)\wedge(-6\gamma+u)
   +6\gamma\wedge(2\gamma+u+w+\delta)\\
 &=(-6\gamma\wedge u-6\gamma\wedge w+u\wedge w)
   +(6\gamma\wedge u+6\gamma\wedge w
     +6\gamma\wedge\delta)\\
 &=u\wedge w+6\gamma\wedge\delta=\eta_R,
\end{align*}
and
\begin{align*}
 (\bigwedge^2T_2)\eta_R
 &=(6\gamma+w)\wedge(-u)
   +6\gamma\wedge(-3\gamma+u+\delta)\\
 &=(-6\gamma\wedge u+u\wedge w)
   +(6\gamma\wedge u+6\gamma\wedge\delta)\\
 &=u\wedge w+6\gamma\wedge\delta=\eta_R.
\end{align*}
The source regular base has monodromy generated by these two transformations,
so the two identities prove the last assertion.
\end{proof}

The line bundle and connection can be written without invoking an existence
label.  On \(\R^4_y\times\C\), impose
\begin{align}
 (y,z)&\sim(y+e_1,e^{12\pi i y^4}z),
 & (y,z)&\sim(y+e_2,e^{2\pi i y^3}z),
 \label{eq:s6cusp-native-line-clutch-one}\\
 (y,z)&\sim(y+e_3,z),
 & (y,z)&\sim(y+e_4,z).
 \label{eq:s6cusp-native-line-clutch-two}
\end{align}
The integer coefficients make the four deck transformations commute: the
only potentially changed factors acquire \(e^{12\pi i}=1\) or
\(e^{2\pi i}=1\).  Denote the quotient line bundle by
\(\mathcal L_{\eta,R}\).  On the covering bundle define
\begin{equation}\label{eq:s6cusp-native-line-connection}
 \nabla_{\eta,R}:=
 \dd-2\pi i\left(6y^1\dd y^4+y^2\dd y^3\right).
\end{equation}

\begin{proposition}[Exact descent, curvature, and native Chern class]
\label{prop:s6cusp-native-line-bundle}
The connection \eqref{eq:s6cusp-native-line-connection} descends through
\eqref{eq:s6cusp-native-line-clutch-one}--
\eqref{eq:s6cusp-native-line-clutch-two}, and
\begin{equation}\label{eq:s6cusp-native-line-curvature}
 F_{\nabla_{\eta,R}}=-2\pi i\,\eta_R.
\end{equation}
With the convention
\(c_1(\mathcal L)=[iF_\nabla/(2\pi)]\),
\begin{equation}\label{eq:s6cusp-native-line-c1}
 c_1(\mathcal L_{\eta,R})=\eta_R,
 \qquad
 \frac12\int_{F_R}c_1(\mathcal L_{\eta,R})^2=6.
\end{equation}
The construction uses the full four-torus and no elliptic quotient.
\end{proposition}

\begin{proof}
Let \(A_R=-2\pi i(6y^1\dd y^4+y^2\dd y^3)\).  Translation by
\(e_1\) changes \(A_R\) by \(-12\pi i\dd y^4\), while the first clutch
factor \(g_1=e^{12\pi i y^4}\) satisfies
\(g_1^{-1}\dd g_1=12\pi i\dd y^4\).  Hence
\[
 g_1^{-1}\dd g_1+A_R(y+e_1)=A_R(y).
\]
Translation by \(e_2\) changes \(A_R\) by
\(-2\pi i\dd y^3\), which is cancelled by
\(g_2^{-1}\dd g_2=2\pi i\dd y^3\).  The remaining two translations
leave \(A_R\) unchanged.  These are exactly the four connection-covariance
equations for the clutch maps, proving descent.  Exterior differentiation
gives
\[
 \dd A_R=-2\pi i
 \left(6\dd y^1\wedge\dd y^4+
 \dd y^2\wedge\dd y^3\right),
\]
which is \eqref{eq:s6cusp-native-line-curvature}.  Substitution into the
stated Chern--Weil convention and then
\eqref{eq:s6cusp-native-eta-square-pfaffian} proves
\eqref{eq:s6cusp-native-line-c1}.
\end{proof}

To retain every coefficient in the magnetic operator, write
\begin{equation}\label{eq:s6cusp-native-physical-two-form}
 \eta_R=\frac12(\Omega_R)_{ij}\dd x^i\wedge\dd x^j,
 \qquad
 \Omega_R=P_R^{-\mathsf T}E_\eta P_R^{-1}.
\end{equation}
Direct multiplication gives the complete matrix
\begin{equation}\label{eq:s6cusp-native-physical-matrix}
 \Omega_R=\frac1{\delta_R}
 \begin{pmatrix}
 0&p_R&0&-6m_R\\
 -p_R&0&6m_R&0\\
 0&-6m_R&0&6T_R\\
 6m_R&0&-6T_R&0
 \end{pmatrix}.
\end{equation}
This identity can be checked without inversion by multiplying the displayed
right side on the left by \(P_R^{\mathsf T}\) and on the right by \(P_R\):
the result is exactly \(E_\eta\).  In particular, all three real parts
\(u_R,a_R,c_R^{\mathrm{re}}\) cancel from the physical curvature matrix as
an exact consequence of the native form \eqref{eq:s6cusp-native-eta}; they
remain present in \(P_R\) and in the marking map.

Reorder the covering coordinates into the two column vectors
\begin{equation}\label{eq:s6cusp-native-XY-blocks}
 X:=\begin{pmatrix}x^1\\x^3\end{pmatrix},
 \qquad
 Y:=\begin{pmatrix}x^2\\x^4\end{pmatrix},
 \end{equation}
and retain the symmetric matrix
\begin{equation}\label{eq:s6cusp-native-K-matrix}
 K_R:=
 \begin{pmatrix}
 p_R&-6m_R\\
 -6m_R&6T_R
 \end{pmatrix}.
\end{equation}
In the order \((X;Y)\), equation
\eqref{eq:s6cusp-native-physical-matrix} is
\begin{equation}\label{eq:s6cusp-native-skew-block}
 \Omega_R=\frac1{\delta_R}
 \begin{pmatrix}0&K_R\\-K_R&0\end{pmatrix}.
\end{equation}
The two eigenvalues of \(K_R\), with their signs retained, are
\begin{align}
 \kappa_{R,+}
 &:=\frac{p_R+6T_R+
 \sqrt{(p_R-6T_R)^2+144m_R^2}}{2}>0,
 \label{eq:s6cusp-kappa-plus}\\
 \kappa_{R,-}
 &:=\frac{p_R+6T_R-
 \sqrt{(p_R-6T_R)^2+144m_R^2}}{2}<0.
 \label{eq:s6cusp-kappa-minus}
\end{align}
Indeed,
\begin{equation}\label{eq:s6cusp-kappa-product}
 \kappa_{R,+}\kappa_{R,-}
 =\det K_R=6p_RT_R-36m_R^2=-6\delta_R<0.
\end{equation}
The exact orthogonal spectral projectors, including their maps and kernels,
are
\begin{align}
 \mathsf Q_{R,+}
 &:=\frac{K_R-\kappa_{R,-}I}
 {\kappa_{R,+}-\kappa_{R,-}},
 &\ker\mathsf Q_{R,+}&=\operatorname{im}\mathsf Q_{R,-},
 \label{eq:s6cusp-native-projector-plus}\\
 \mathsf Q_{R,-}
 &:=\frac{\kappa_{R,+}I-K_R}
 {\kappa_{R,+}-\kappa_{R,-}},
 &\ker\mathsf Q_{R,-}&=\operatorname{im}\mathsf Q_{R,+}.
 \label{eq:s6cusp-native-projector-minus}
\end{align}
They satisfy
\begin{equation}\label{eq:s6cusp-native-projector-identities}
 \mathsf Q_{R,+}+\mathsf Q_{R,-}=I,
 \quad
 \mathsf Q_{R,\sigma}\mathsf Q_{R,\rho}
 =\delta_{\sigma\rho}\mathsf Q_{R,\sigma},
 \quad
 K_R\mathsf Q_{R,\sigma}
 =\kappa_{R,\sigma}\mathsf Q_{R,\sigma}.
\end{equation}
Thus the orthogonal coordinate map
\begin{equation}\label{eq:s6cusp-native-orthogonal-map}
 \mathscr O_R(X,Y):=
 \bigl(
 \mathsf Q_{R,+}X,\mathsf Q_{R,+}Y;
 \mathsf Q_{R,-}X,-\mathsf Q_{R,-}Y
 \bigr)
\end{equation}
has inverse
\begin{equation}\label{eq:s6cusp-native-orthogonal-inverse}
 (X_+,Y_+;X_-,Y_-)
 \longmapsto(X_++X_-,Y_+-Y_-).
\end{equation}
It sends the real magnetic two-form \(2\pi\eta_R\) to the orthogonal sum
of two positively oriented constant fields with exact frequencies
\begin{align}
 \mathcal B_{R,+}
 &:=\frac{2\pi\kappa_{R,+}}{\delta_R}
 =\frac{\pi}{\delta_R}
 \left[p_R+6T_R+
 \sqrt{(p_R-6T_R)^2+144m_R^2}\right],
 \label{eq:s6cusp-native-frequency-plus}\\
 \mathcal B_{R,-}
 &:=-\frac{2\pi\kappa_{R,-}}{\delta_R}
 =\frac{\pi}{\delta_R}
 \left[-p_R-6T_R+
 \sqrt{(p_R-6T_R)^2+144m_R^2}\right].
 \label{eq:s6cusp-native-frequency-minus}
\end{align}

\begin{theorem}[Complete native degree-six four-dimensional Landau spectrum]
\label{thm:s6cusp-native-four-dimensional-landau}
Let \(H_{\eta,R}\) be the Friedrichs realization of
\begin{equation}\label{eq:s6cusp-native-magnetic-Hamiltonian}
 H_{\eta,R}:=\nabla_{\eta,R}^*\nabla_{\eta,R}
\end{equation}
on \(L^2(F_R,\mathcal L_{\eta,R})\), with the Euclidean covering metric
\(g_R\).  Its joint Landau levels, counted before coincident numerical
energies are combined, are
\begin{equation}\label{eq:s6cusp-native-landau-spectrum}
 E_{R;n_+,n_-}
 =(2n_++1)\mathcal B_{R,+}
 +(2n_-+1)\mathcal B_{R,-},
 \qquad n_+,n_-\in\Z_{\ge0}.
\end{equation}
Every joint eigenspace in
\eqref{eq:s6cusp-native-landau-spectrum} has complex dimension six.  If
two different index pairs give the same numerical energy, their six-
dimensional joint eigenspaces add orthogonally in that energy eigenspace.
The ground space has complex dimension six, and the exact excitation gap is
\begin{equation}\label{eq:s6cusp-native-landau-gap}
 \operatorname{gap}(H_{\eta,R})
 =2\min\{\mathcal B_{R,+},\mathcal B_{R,-}\}.
\end{equation}
The product of the two retained frequencies is
\begin{align}
 \mathcal B_{R,+}\mathcal B_{R,-}
 &=\frac{24\pi^2}{\delta_R}
 \label{eq:s6cusp-native-frequency-product}\\
 &=-\frac{24\pi^2}{
 \operatorname{Im}\tau_R
 \left(
 \operatorname{Im}\beta_R-
 \dfrac{6(\operatorname{Im}\mu_R)^2}
 {\operatorname{Im}\tau_R}
 \right)}.
 \label{eq:s6cusp-native-frequency-unchanged-expression}
\end{align}
Along the source radial cusp,
\begin{align}
 \mathcal B_{R,+}
 &=\frac{12\pi}{R}\left[1+O(R^{-1})\right],
 &
 \mathcal B_{R,-}
 &=\frac{2\pi}{R}\left[1+O(R^{-1})\right].
 \label{eq:s6cusp-native-frequency-asymptotics}
\end{align}
Consequently there is \(R_3\) such that, for \(R\ge R_3\),
\begin{equation}\label{eq:s6cusp-native-gap-exact-eventual}
 \operatorname{gap}(H_{\eta,R})
 =\frac{2\pi}{\delta_R}
 \left[
 \sqrt{(p_R-6T_R)^2+144m_R^2}-p_R-6T_R
 \right]
 =\frac{4\pi}{R}\left[1+O(R^{-1})\right].
\end{equation}
\end{theorem}

\begin{proof}
The closed quadratic form
\(q_{\eta,R}[s]=\lVert\nabla_{\eta,R}s\rVert_{L^2}^2\) has domain
\(H^1(F_R,\mathcal L_{\eta,R})\).  In a finite unitary trivializing cover,
its graph norm is equivalent to the ordinary first Sobolev norm because the
connection one-forms are smooth and bounded.  The Rellich inclusion on the
compact torus is compact, so the associated Friedrichs operator is
self-adjoint with compact resolvent and smooth eigensections.

Apply the orthogonal map \(\mathscr O_R\) on the universal cover.  Choose a
unit vector in each one-dimensional image
\(\operatorname{im}\mathsf Q_{R,\sigma}\); changing its sign changes both
coordinates in that block and leaves every operator below unchanged.  Let
\(\Pi_{X_\sigma},\Pi_{Y_\sigma}\) be the corresponding self-adjoint
covariant momenta.  Equations
\eqref{eq:s6cusp-native-skew-block}--
\eqref{eq:s6cusp-native-frequency-minus} give the complete commutator table
\begin{equation}\label{eq:s6cusp-native-magnetic-commutators}
 [\Pi_{X_\sigma},\Pi_{Y_\rho}]
 =i\delta_{\sigma\rho}\mathcal B_{R,\sigma},
 \qquad
 [\Pi_{X_\sigma},\Pi_{X_\rho}]
 =[\Pi_{Y_\sigma},\Pi_{Y_\rho}]=0.
\end{equation}
Define on smooth sections
\begin{equation}\label{eq:s6cusp-native-ladders}
 \mathfrak a_\sigma
 :=\frac{\Pi_{X_\sigma}+i\Pi_{Y_\sigma}}
 {\sqrt{2\mathcal B_{R,\sigma}}},
 \qquad
 \mathfrak a_\sigma^*
 :=\frac{\Pi_{X_\sigma}-i\Pi_{Y_\sigma}}
 {\sqrt{2\mathcal B_{R,\sigma}}}.
\end{equation}
Direct multiplication, with both frequencies retained, gives
\begin{align}
 [\mathfrak a_\sigma,\mathfrak a_\rho^*]
 &=\delta_{\sigma\rho}I,
 &
 [\mathfrak a_\sigma,\mathfrak a_\rho]&=0,
 \label{eq:s6cusp-native-ladder-commutators}\\
 H_{\eta,R}
 &=\mathcal B_{R,+}
 (2\mathfrak a_+^*\mathfrak a_++I)
 +\mathcal B_{R,-}
 (2\mathfrak a_-^*\mathfrak a_-+I).
 \label{eq:s6cusp-native-oscillator-factorization}
\end{align}

Put \(\mathfrak N_\sigma=\mathfrak a_\sigma^*\mathfrak a_\sigma\).
The commutators show that the two number operators commute and that
lowering a nonzero eigenvector decreases the corresponding eigenvalue by
one.  If an eigenvalue \(\nu\) of one number operator were nonintegral,
successive lowering through \(\lfloor\nu\rfloor+1\) steps would produce a
nonzero vector with negative eigenvalue, contradicting
\(\mathfrak N_\sigma\ge0\).  For \(n\ge1\), the maps
\begin{equation}\label{eq:s6cusp-native-level-isomorphisms}
 n^{-1/2}\mathfrak a_\sigma:
 \ker(\mathfrak N_\sigma-nI)
 \underset{n^{-1/2}\mathfrak a_\sigma^*}{\rightleftarrows}
 \ker(\mathfrak N_\sigma-(n-1)I)
\end{equation}
are mutually inverse while preserving the other number eigenvalue.  Thus
all joint eigenspaces have the dimension of
\begin{equation}\label{eq:s6cusp-native-joint-ground}
 \mathcal K_R:=\ker\mathfrak a_+\cap\ker\mathfrak a_-.
\end{equation}

For completeness, equip the underlying real torus with the constant complex
structure for which the two oriented planes in
\eqref{eq:s6cusp-native-orthogonal-map} are complex lines.  The form
\(2\pi\eta_R\) is positive for this complex structure, and the two equations
in \eqref{eq:s6cusp-native-joint-ground} are exactly the covariant
\(\bar\partial\)-equations.  Hence
\(\mathcal K_R\cong H^0(F_R,\mathcal L_{\eta,R})\) for this retained
positive magnetic complex structure.  The canonical bundle of a complex
torus is trivial.  The Bochner--Kodaira identity for the positive curvature
annihilates the higher cohomology, while Hirzebruch--Riemann--Roch has
Todd class one and gives
\[
 \dim_\C\mathcal K_R
 =\chi(\mathcal L_{\eta,R})
 =\frac12\int_{F_R}c_1(\mathcal L_{\eta,R})^2=6.
\]
This is the dimension-six theta-space calculation; the classical
polarized-torus decomposition and its theta basis are recorded at source
level in \cite{Tyurin1999}.  Raising from the six ground vectors and using
\eqref{eq:s6cusp-native-level-isomorphisms} proves the joint spectrum and
every multiplicity in \eqref{eq:s6cusp-native-landau-spectrum}.

The smallest increase from \((n_+,n_-)=(0,0)\) is twice the smaller
frequency, proving \eqref{eq:s6cusp-native-landau-gap}.  Equation
\eqref{eq:s6cusp-kappa-product} gives
\[
 \mathcal B_{R,+}\mathcal B_{R,-}
  =\frac{(2\pi)^2(-\kappa_{R,+}\kappa_{R,-})}{\delta_R^2}
  =\frac{24\pi^2}{\delta_R}.
\]
Substitution of the unchanged identity
\eqref{eq:s6cusp-unchanged-expression-delta} gives
\eqref{eq:s6cusp-native-frequency-unchanged-expression} without deleting
its parenthesis.

Finally, equations \eqref{eq:s6cusp-imaginary-identities} give
\[
 \frac{p_R}{R}\longrightarrow-1,\qquad
 \frac{T_R}{R}\longrightarrow1,\qquad
 \frac{m_R}{R}\longrightarrow0,\qquad
 \frac{\delta_R}{R^2}\longrightarrow1.
\]
Substitution into \eqref{eq:s6cusp-kappa-plus}--
\eqref{eq:s6cusp-native-frequency-minus} proves
\eqref{eq:s6cusp-native-frequency-asymptotics}.  Their ratio tends to six,
so \(\mathcal B_{R,-}<\mathcal B_{R,+}\) for all sufficiently large
\(R\).  Doubling the exact right side of
\eqref{eq:s6cusp-native-frequency-minus} then gives
\eqref{eq:s6cusp-native-gap-exact-eventual}.
\end{proof}

The source complex structure makes the same integral class \(\eta_R\) of
type \((1,1)\) with indefinite Hermitian form; the magnetic complex
structure used only in the multiplicity proof makes its curvature positive.
The identity of the underlying marked smooth torus and the identity of the
smooth clutch bundle are the exact maps relating those two presentations.
No equality of their complex structures is used.  The spectrum
\eqref{eq:s6cusp-native-landau-spectrum} is that of the smooth unitary
connection \eqref{eq:s6cusp-native-line-connection} and the auxiliary
Euclidean metric \(g_R\), not of a Yang--Mills transfer Hamiltonian.

\subsection{Exact heat trace and the threshold-correct fixed-energy count}
\label{subsec:s6cusp-native-heat-count}

Retain the absolute ground energy
\begin{equation}\label{eq:s6cusp-native-ground-energy}
 E_{R;0,0}:=\mathcal B_{R,+}+\mathcal B_{R,-}
 =\frac{2\pi}{\delta_R}
 \sqrt{(p_R-6T_R)^2+144m_R^2}.
\end{equation}
The adjective ``ground'' in this definition is justified by the quadratic-
form identity \eqref{eq:s6cusp-native-oscillator-factorization}: for every
\(s\) in the form domain,
\begin{align}
 \left\langle s,
 (H_{\eta,R}-E_{R;0,0})s\right\rangle
 &=2\mathcal B_{R,+}\lVert\mathfrak a_+s\rVert^2
   +2\mathcal B_{R,-}\lVert\mathfrak a_-s\rVert^2
 \label{eq:s6cusp-native-ground-lower-bound}\\
 &\ge0.
\end{align}
Both coefficients are strictly positive.  Indeed,
\(\delta_R>0\) and
\(\Xi_R^2-S_R^2=24\delta_R\), where
\begin{equation}\label{eq:s6cusp-Xi-S-definitions}
 \Xi_R:=\sqrt{(p_R-6T_R)^2+144m_R^2},
 \qquad S_R:=p_R+6T_R,
\end{equation}
imply \(\Xi_R>|S_R|\).  Equations
\eqref{eq:s6cusp-native-frequency-plus}--
\eqref{eq:s6cusp-native-frequency-minus} therefore give
\(\mathcal B_{R,+},\mathcal B_{R,-}>0\).  The six-dimensional simultaneous
kernel \eqref{eq:s6cusp-native-joint-ground} proves that the lower bound is
attained with complex multiplicity six.  Thus topology supplies six states
\emph{at} \(E_{R;0,0}\); it supplies no state below that number.

\begin{theorem}[Exact native heat trace]
\label{thm:s6cusp-native-heat-trace}
For every \(t>0\),
\begin{equation}\label{eq:s6cusp-native-heat-trace}
 \operatorname{Tr}(e^{-tH_{\eta,R}})
 =\frac{3}{
  2\sinh(t\mathcal B_{R,+})
   \sinh(t\mathcal B_{R,-})}.
\end{equation}
At every fixed \(t>0\), the source radial cusp has the expansion
\begin{equation}\label{eq:s6cusp-native-heat-trace-cusp}
 \operatorname{Tr}(e^{-tH_{\eta,R}})
 =\frac{\delta_R}{16\pi^2t^2}
 \left[1+O_t(R^{-2})\right].
\end{equation}
\end{theorem}

\begin{proof}
The complete spectrum and its six-dimensional joint multiplicities give
\begin{align*}
 \operatorname{Tr}(e^{-tH_{\eta,R}})
 &=6\sum_{n_+,n_-\ge0}
 e^{-t[(2n_++1)\mathcal B_{R,+}
       +(2n_-+1)\mathcal B_{R,-}]}\\
 &=6\prod_{\sigma\in\{+,-\}}
 \frac{e^{-t\mathcal B_{R,\sigma}}}
 {1-e^{-2t\mathcal B_{R,\sigma}}}\\
 &=\frac{3}{2\sinh(t\mathcal B_{R,+})
                  \sinh(t\mathcal B_{R,-})},
\end{align*}
which proves \eqref{eq:s6cusp-native-heat-trace}.  The exact product
\eqref{eq:s6cusp-native-frequency-product} and
\(\mathcal B_{R,+},\mathcal B_{R,-}=O(R^{-1})\) give
\begin{align*}
 \sinh(t\mathcal B_{R,+})
 \sinh(t\mathcal B_{R,-})
 &=t^2\mathcal B_{R,+}\mathcal B_{R,-}
 \left[1+O_t(\mathcal B_{R,+}^2+
                  \mathcal B_{R,-}^2)\right]\\
 &=\frac{24\pi^2t^2}{\delta_R}
 \left[1+O_t(R^{-2})\right].
\end{align*}
Substitution into the exact trace proves
\eqref{eq:s6cusp-native-heat-trace-cusp}.
\end{proof}


For \(E>0\), define the absolute spectral counting function
\begin{equation}\label{eq:s6cusp-native-count-definition}
 N_{\eta,R}(E):=
 \operatorname{Tr}\mathbf1_{(-\infty,E]}(H_{\eta,R}).
\end{equation}
The exact lattice of admissible indices is
\begin{equation}\label{eq:s6cusp-native-count-index-set}
 \mathscr I_R(E):=
 \left\{(n_+,n_-)\in\mathbb Z_{\ge0}^2:
  2\mathcal B_{R,+}n_++2\mathcal B_{R,-}n_-
  \le E-E_{R;0,0}\right\}.
\end{equation}
It is empty when \(E<E_{R;0,0}\).  Consequently
\begin{equation}\label{eq:s6cusp-native-count-cardinality}
 N_{\eta,R}(E)=6\,\#\mathscr I_R(E).
\end{equation}
Writing
\begin{equation}\label{eq:s6cusp-native-count-L}
 L_R(E):=E-E_{R;0,0}
\end{equation}
without deleting its sign, the exact floor formula is
\begin{equation}\label{eq:s6cusp-native-count-piecewise}
 N_{\eta,R}(E)=
 \begin{cases}
 0,&E<E_{R;0,0},\\[0.4em]
 6\displaystyle\sum_{n_+=0}^{
   \left\lfloor L_R(E)/(2\mathcal B_{R,+})\right\rfloor}
 \left[
  1+\left\lfloor
  \dfrac{L_R(E)-2\mathcal B_{R,+}n_+}
        {2\mathcal B_{R,-}}
  \right\rfloor
 \right],&E\ge E_{R;0,0}.
 \end{cases}
\end{equation}
In particular,
\begin{equation}\label{eq:s6cusp-native-count-threshold-values}
 \operatorname{Tr}\mathbf1_{(-\infty,E_{R;0,0})}(H_{\eta,R})=0,
 \qquad
 N_{\eta,R}(E_{R;0,0})=6.
\end{equation}
Replacing \(L_R(E)\) by
\((E-E_{R;0,0})_+\) inside the second line while using that line for every
\(E>0\) would identify the two distinct cases in
\eqref{eq:s6cusp-native-count-threshold-values}: the zero value of the
positive part would retain the index pair \((0,0)\) even below the spectrum.

\begin{theorem}[Fixed-energy accumulation with the threshold retained]
\label{thm:s6cusp-native-fixed-energy-count}
For every fixed \(E>0\), as \(R\to\infty\),
\begin{align}
 N_{\eta,R}(E)
 &=\frac{\delta_R}{32\pi^2}
   \bigl(E-E_{R;0,0}\bigr)^2+O_E(R)
 \label{eq:s6cusp-native-fixed-energy-count-ground}\\
 &=\frac{E^2\delta_R}{32\pi^2}+O_E(R).
 \label{eq:s6cusp-native-fixed-energy-count-leading}
\end{align}
\end{theorem}

\begin{proof}
Equation \eqref{eq:s6cusp-native-ground-energy} and
\eqref{eq:s6cusp-native-frequency-asymptotics} give
\(E_{R;0,0}=14\pi R^{-1}+O(R^{-2})\).  Hence every fixed \(E>0\)
eventually satisfies \(E>E_{R;0,0}\), and the second branch of
\eqref{eq:s6cusp-native-count-piecewise} applies.  It counts the integer
points in the closed right triangle
\begin{equation}\label{eq:s6cusp-native-count-triangle}
 \left\{(x,y)\in\mathbb R_{\ge0}^2:
 2\mathcal B_{R,+}x+2\mathcal B_{R,-}y
 \le E-E_{R;0,0}\right\}.
\end{equation}
Its area, with both unequal intercepts retained, is
\begin{equation}\label{eq:s6cusp-native-count-triangle-area}
 \frac{(E-E_{R;0,0})^2}
 {8\mathcal B_{R,+}\mathcal B_{R,-}}.
\end{equation}
Attach the half-open unit square
\([n_+,n_++1)\times[n_-,n_-+1)\) to each counted lattice point.
The symmetric difference between their union and the triangle is contained
in a fixed-width neighborhood of the three boundary segments.  Its area is
bounded by a universal constant times the perimeter plus one.  The two axis
intercepts are
\begin{equation*}
 \frac{E-E_{R;0,0}}{2\mathcal B_{R,+}},
 \qquad
 \frac{E-E_{R;0,0}}{2\mathcal B_{R,-}},
\end{equation*}
and are \(O_E(R)\), so the lattice-point discrepancy is \(O_E(R)\).
Multiplication by six and the exact identity
\(\mathcal B_{R,+}\mathcal B_{R,-}=24\pi^2/\delta_R\)
prove \eqref{eq:s6cusp-native-fixed-energy-count-ground}.  Finally,
\begin{align*}
 \delta_R\left[(E-E_{R;0,0})^2-E^2\right]
 &=-2E\delta_RE_{R;0,0}
   +\delta_RE_{R;0,0}^2\\
 &=O_E(R),
\end{align*}
because \(\delta_R=R^2+O(R)\) and
\(E_{R;0,0}=14\pi R^{-1}+O(R^{-2})\).  This proves
\eqref{eq:s6cusp-native-fixed-energy-count-leading}.
\end{proof}

\subsection{The native topological split \(SU(2)\) saddle}
\label{subsec:s6cusp-native-split-su2}

Form the rank-two bundle
\begin{equation}\label{eq:s6cusp-native-split-bundle}
 \mathcal E_{\eta,R}:=
 \mathcal L_{\eta,R}\oplus\mathcal L_{\eta,R}^{-1}.
\end{equation}
The determinant trivialization is the evaluation pairing between the two
summands.  Retaining both Chern roots gives
\begin{align}
 c(\mathcal E_{\eta,R})
 &=(1+\eta_R)(1-\eta_R)=1-\eta_R^2,
 \label{eq:s6cusp-native-split-total-chern}\\
 c_1(\mathcal E_{\eta,R})&=0,
 &c_2(\mathcal E_{\eta,R})&=-\eta_R^2,
 &\int_{F_R}c_2(\mathcal E_{\eta,R})&=-12.
 \label{eq:s6cusp-native-split-c2}
\end{align}
Thus this split \(SU(2)\) bundle is not the trivial rank-two bundle: a bundle
isomorphism would pull back the zero second Chern class of the latter to the
nonzero class in \eqref{eq:s6cusp-native-split-c2}.

Put
\begin{equation}\label{eq:s6cusp-native-Hmag}
 H_{\mathrm{mag}}:=\operatorname{diag}(i,-i),
 \qquad
 \langle X,Y\rangle_{\mathfrak{su}(2)}
 :=-\frac12\operatorname{Tr}(XY),
\end{equation}
so that \(\lVert H_{\mathrm{mag}}\rVert=1\).  The direct-sum connection
\begin{equation}\label{eq:s6cusp-native-split-connection}
 A_{\eta,R}:=
 \nabla_{\eta,R}\oplus\nabla_{\eta,R}^{-1}
\end{equation}
has the exact curvature
\begin{equation}\label{eq:s6cusp-native-split-curvature}
 F_{A_{\eta,R}}=-2\pi\eta_RH_{\mathrm{mag}}.
\end{equation}
All coefficients of \(F_{A_{\eta,R}}\) in the covering coordinates are
constant, and both connection summands commute with
\(H_{\mathrm{mag}}\).  Therefore
\begin{equation}\label{eq:s6cusp-native-split-YM-equations}
 \dd_{A_{\eta,R}}F_{A_{\eta,R}}=0,
 \qquad
 \dd_{A_{\eta,R}}^*F_{A_{\eta,R}}=0.
\end{equation}
The first identity is also the Bianchi identity; the second follows directly
from the constant canonical coefficients and the flat covering metric.

Use throughout this subsection the action
\begin{equation}\label{eq:s6cusp-native-YM-action}
 \operatorname{YM}_R(A):=
 \frac1{2g_{\mathrm{YM}}^2}
 \int_{F_R}\lVert F_A\rVert^2\operatorname{vol}_{g_R}.
\end{equation}

\begin{theorem}[Exact action and exact topological excess]
\label{thm:s6cusp-native-split-action}
The Yang--Mills critical connection \(A_{\eta,R}\) satisfies
\begin{align}
 \operatorname{YM}_R(A_{\eta,R})
 &=\frac{\delta_R}{2g_{\mathrm{YM}}^2}
   \left(\mathcal B_{R,+}^2+\mathcal B_{R,-}^2\right)
 \label{eq:s6cusp-native-action-frequencies}\\
 &=\frac{2\pi^2}{g_{\mathrm{YM}}^2\delta_R}
   \left(p_R^2+36T_R^2+72m_R^2\right)
 \label{eq:s6cusp-native-action-coordinates}\\
 &=\frac{24\pi^2}{g_{\mathrm{YM}}^2}
   +\frac{2\pi^2(p_R+6T_R)^2}
   {g_{\mathrm{YM}}^2\delta_R}.
 \label{eq:s6cusp-native-action-excess}
\end{align}
Along the source radial cusp,
\begin{equation}\label{eq:s6cusp-native-action-cusp}
 \operatorname{YM}_R(A_{\eta,R})
 =\frac{74\pi^2}{g_{\mathrm{YM}}^2}
  +\frac{70\pi^2d_0}{g_{\mathrm{YM}}^2R}
  +O(R^{-2}).
\end{equation}
\end{theorem}

\begin{proof}
In the positive canonical planes, the two signed coefficients of \(\eta_R\)
have absolute values
\(\mathcal B_{R,+}/(2\pi)\) and
\(\mathcal B_{R,-}/(2\pi)\).  Equations
\eqref{eq:s6cusp-native-Hmag} and
\eqref{eq:s6cusp-native-split-curvature} therefore give
\begin{equation*}
 \lVert F_{A_{\eta,R}}\rVert^2
 =\mathcal B_{R,+}^2+\mathcal B_{R,-}^2.
\end{equation*}
The marked Euclidean volume is
\(|\det P_R|=\delta_R\), proving
\eqref{eq:s6cusp-native-action-frequencies}.  Since the signed canonical
coefficients are the eigenvalues of \(K_R/\delta_R\),
\begin{equation*}
 \kappa_{R,+}^2+\kappa_{R,-}^2
 =\operatorname{Tr}(K_R^2)
 =p_R^2+36T_R^2+72m_R^2.
\end{equation*}
Substitution proves \eqref{eq:s6cusp-native-action-coordinates}.  Direct
expansion, with \(p_R,T_R,m_R\) retained, gives
\begin{equation}\label{eq:s6cusp-native-action-coordinate-decomposition}
 p_R^2+36T_R^2+72m_R^2
 =12(6m_R^2-p_RT_R)+(p_R+6T_R)^2
 =12\delta_R+(p_R+6T_R)^2.
\end{equation}
This proves \eqref{eq:s6cusp-native-action-excess}.

For the radial expansion, equations
\eqref{eq:s6cusp-imaginary-identities} and
\eqref{eq:s6cusp-delta-two-term-expansion} retain
\begin{align*}
 p_R+6T_R
 &=5R+(d_0+5h_0)+O(e^{-2\pi R}),\\
 \delta_R
 &=R^2+(2h_0-d_0)R
   +(h_0^2-d_0h_0+6m_0^2)+O(Re^{-2\pi R}).
\end{align*}
Long division without deleting either displayed polynomial yields
\begin{equation*}
 \frac{(p_R+6T_R)^2}{\delta_R}
 =25+\frac{35d_0}{R}+O(R^{-2}).
\end{equation*}
Inserting this in \eqref{eq:s6cusp-native-action-excess} proves
\eqref{eq:s6cusp-native-action-cusp}.
\end{proof}

\subsection{Complete Coulomb-reduced Jacobi index of the native saddle}
\label{subsec:s6cusp-native-jacobi-index}

Let
\begin{equation}\label{eq:s6cusp-native-adjoint-root-splitting}
 \operatorname{ad}(\mathcal E_{\eta,R})_{\mathbb C}
 =\underline{\mathbb C}H_{\mathrm{mag}}
 \oplus\mathcal L_{\eta,R}^{2}E_+
 \oplus\mathcal L_{\eta,R}^{-2}E_-,
 \qquad
 [H_{\mathrm{mag}},E_\pm]=\pm2iE_\pm.
\end{equation}
The root-line Chern matrix is \(2E_\eta\), and therefore
\begin{equation}\label{eq:s6cusp-native-root-pfaffian}
 \left|\operatorname{Pf}(2E_\eta)\right|
 =4\left|\operatorname{Pf}(E_\eta)\right|=24.
\end{equation}
Applying the proved oscillator construction of
\cref{thm:s6cusp-native-four-dimensional-landau} to
\(\mathcal L_{\eta,R}^{2}\) gives scalar root levels
\begin{equation}\label{eq:s6cusp-native-root-scalar-levels}
 \lambda^{(0)}_{R;n_+,n_-}
 =2(2n_++1)\mathcal B_{R,+}
  +2(2n_-+1)\mathcal B_{R,-},
\end{equation}
each with complex multiplicity twenty-four.

On the complexified one-form Hilbert space define the full gauge-fixed
Jacobi operator with domain \(H^2\) by
\begin{equation}\label{eq:s6cusp-native-full-jacobi}
 (\mathcal J^{\mathrm{gf}}_{\eta,R}a)_\nu
 :=-D_{A_{\eta,R}}^\mu D_{A_{\eta,R},\mu}a_\nu
   -2[F_{A_{\eta,R},\nu\mu},a_\mu].
\end{equation}
This is a symmetric elliptic Laplace-type operator on a smooth bundle over
the compact torus, so its \(H^2\) realization is self-adjoint and has compact
resolvent.  Define the closed Coulomb Hilbert subspace by
\begin{equation}\label{eq:s6cusp-native-Coulomb-Hilbert-space}
 \mathscr C_{\eta,R}:=
 \left(\overline{\operatorname{im}
  (\dd_{A_{\eta,R}}:H^1(\operatorname{ad}\mathcal E_{\eta,R})
   \longrightarrow L^2(T^*F_R\otimes\operatorname{ad}\mathcal E_{\eta,R}))}
 \right)^\perp.
\end{equation}
For an \(H^1\) one-form, membership is equivalent to the weak equation
\(\dd_{A_{\eta,R}}^*a=0\).

The Coulomb space reduces \(\mathcal J^{\mathrm{gf}}_{\eta,R}\).  To prove
this, let \(\mathcal L^{\mathrm{YM}}_{\eta,R}\) denote the ungauged
Yang--Mills Hessian differential at the critical connection.  Gauge
covariance of
\(A\mapsto\dd_A^*F_A\) gives, after differentiating the path
\(\exp(t\phi)\cdot A_{\eta,R}\) at \(t=0\),
\begin{equation}\label{eq:s6cusp-native-gauge-orbit-hessian-zero}
 \mathcal L^{\mathrm{YM}}_{\eta,R}\dd_{A_{\eta,R}}\phi=0.
\end{equation}
The gauge-fixed operator is
\(\mathcal J^{\mathrm{gf}}_{\eta,R}
=\mathcal L^{\mathrm{YM}}_{\eta,R}
 +\dd_{A_{\eta,R}}\dd_{A_{\eta,R}}^*\).
Hence, on smooth zero-forms,
\begin{equation}\label{eq:s6cusp-native-jacobi-gradient-intertwiner}
 \mathcal J^{\mathrm{gf}}_{\eta,R}\dd_{A_{\eta,R}}
 =\dd_{A_{\eta,R}}\Delta_{A_{\eta,R},0},
 \qquad
 \Delta_{A_{\eta,R},0}:=
 \dd_{A_{\eta,R}}^*\dd_{A_{\eta,R}}.
\end{equation}
For smooth \(a\in\mathscr C_{\eta,R}\) and smooth \(\phi\), self-adjointness
and \eqref{eq:s6cusp-native-jacobi-gradient-intertwiner} give
\begin{align*}
 \left\langle \dd_{A_{\eta,R}}\phi,
 \mathcal J^{\mathrm{gf}}_{\eta,R}a\right\rangle
 &=\left\langle
 \mathcal J^{\mathrm{gf}}_{\eta,R}\dd_{A_{\eta,R}}\phi,a
 \right\rangle\\
 &=\left\langle
 \dd_{A_{\eta,R}}\Delta_{A_{\eta,R},0}\phi,a
 \right\rangle=0.
\end{align*}
Density and closedness prove invariance of \(\mathscr C_{\eta,R}\); the
orthogonal complement is invariant as well.  The restricted operator
\begin{equation}\label{eq:s6cusp-native-Coulomb-Jacobi}
 \mathcal J^{\mathrm C}_{\eta,R}:=
 \mathcal J^{\mathrm{gf}}_{\eta,R}\big|_{
 \operatorname{Dom}(\mathcal J^{\mathrm{gf}}_{\eta,R})
 \cap\mathscr C_{\eta,R}}
\end{equation}
is therefore self-adjoint with discrete spectrum.

\begin{theorem}[Three negative bands and exact real Morse index]
\label{thm:s6cusp-native-jacobi-index}
There is \(R_4\) such that every \(R\ge R_4\) has
\begin{equation}\label{eq:s6cusp-native-ratio-five-seven}
 5<\frac{\mathcal B_{R,+}}{\mathcal B_{R,-}}<7.
\end{equation}
For those \(R\), the complete negative spectrum of
\(\mathcal J^{\mathrm C}_{\eta,R}\) consists of the three distinct
eigenvalues
\begin{align}
 \nu_{R,n}
 &:=2\left[(2n+1)\mathcal B_{R,-}-\mathcal B_{R,+}\right]
 \label{eq:s6cusp-native-negative-bands}\\
 &=\frac{4\pi}{\delta_R}
 \left[n\Xi_R-(n+1)S_R\right],
 \qquad n=0,1,2.
 \label{eq:s6cusp-native-negative-bands-coordinates}
\end{align}
Every one of these eigenspaces has real dimension forty-eight.  Consequently
\begin{equation}\label{eq:s6cusp-native-Morse-index}
 \operatorname{index}_{\mathbb R}
 (\mathcal J^{\mathrm C}_{\eta,R})=3\cdot48=144.
\end{equation}
Their leading cusp terms, together with the next positive member of the same
branch, are
\begin{align}
 \nu_{R,0}&=-\frac{20\pi}{R}+O(R^{-2}),&
 \nu_{R,1}&=-\frac{12\pi}{R}+O(R^{-2}),\notag\\
 \nu_{R,2}&=-\frac{4\pi}{R}+O(R^{-2}),&
 \nu_{R,3}&= \frac{4\pi}{R}+O(R^{-2}).
 \label{eq:s6cusp-native-negative-band-cusp}
\end{align}
\end{theorem}

\begin{proof}
The frequency ratio tends to six by
\eqref{eq:s6cusp-native-frequency-asymptotics}; its strict containment in
\((5,7)\) follows for all \(R\) beyond one fixed \(R_4\).

Use the exact canonical coordinates supplied by
\eqref{eq:s6cusp-native-orthogonal-map}.  On the \(E_+\) root line, the
spin-curvature part of \eqref{eq:s6cusp-native-full-jacobi} acts on the
\((\dd X_+,\dd Y_+)\) coefficients by
\begin{equation}\label{eq:s6cusp-native-strong-spin-matrix}
 \mathsf S_{R,+}
 =4\mathcal B_{R,+}
 \begin{pmatrix}0&i\\-i&0\end{pmatrix},
\end{equation}
and on the \((\dd X_-,\dd Y_-)\) coefficients by the same matrix with
\(\mathcal B_{R,+}\) replaced by \(\mathcal B_{R,-}\).  Direct
multiplication gives
\begin{align*}
 \mathsf S_{R,+}\binom{1}{i}
 &=-4\mathcal B_{R,+}\binom{1}{i},&
 \mathsf S_{R,+}\binom{1}{-i}
 &=4\mathcal B_{R,+}\binom{1}{-i}.
\end{align*}
Thus the four constant circular one-form polarizations shift the scalar
levels \eqref{eq:s6cusp-native-root-scalar-levels} by
\(-4\mathcal B_{R,+},+4\mathcal B_{R,+},
-4\mathcal B_{R,-},+4\mathcal B_{R,-}\), respectively.

For the negative strong-plane polarization the complete level is
\begin{equation}\label{eq:s6cusp-native-strong-negative-full-level}
 2\left[(2n_+-1)\mathcal B_{R,+}
       +(2n_-+1)\mathcal B_{R,-}\right].
\end{equation}
When \(n_+\ge1\), both coefficients in brackets are positive.  When
\(n_+=0\), negativity is equivalent to
\begin{equation*}
 2n_-+1<\frac{\mathcal B_{R,+}}{\mathcal B_{R,-}}.
\end{equation*}
Under \eqref{eq:s6cusp-native-ratio-five-seven}, the odd integers satisfying
this inequality are exactly \(1,3,5\), so
\(n_-=0,1,2\).  The next odd integer is seven, and the strict upper bound in
\eqref{eq:s6cusp-native-ratio-five-seven} makes the \(n_-=3\) level
positive.  The weak-plane negative polarization has minimum
\begin{equation*}
 2(\mathcal B_{R,+}-\mathcal B_{R,-})>0,
\end{equation*}
the two positive spin shifts are positive, and the Cartan sector is the
nonnegative ordinary Hodge Laplacian.  This exhausts the Cartan summand, both
root summands, and all four one-form polarizations.

It remains to show that restricting to Coulomb gauge removes none of the
negative spectrum.  On \(\mathcal L_{\eta,R}^{2}\), put
\begin{equation*}
 \mathfrak b_+:=
 \frac{\Pi_{X_+}+i\Pi_{Y_+}}
 {\sqrt{4\mathcal B_{R,+}}}.
\end{equation*}
Every \(n_+=0\) scalar state \(\phi_{0,n_-}\) is killed by
\(\mathfrak b_+\).  Its negative-polarization one-form obeys the exact
divergence identity
\begin{align}
 \dd_{A_{\eta,R}}^*
 \left[\phi_{0,n_-}(\dd X_++i\dd Y_+)E_+\right]
 &=-i\sqrt{4\mathcal B_{R,+}}
   (\mathfrak b_+\phi_{0,n_-})E_+\notag\\
 &=0.
 \label{eq:s6cusp-native-negative-Coulomb-equation}
\end{align}
The Hermitian-adjoint \(E_-\) modes satisfy the conjugate equation.  Hence
the complete negative spectral subspace of the full operator lies in the
reducing Coulomb subspace.

Equation \eqref{eq:s6cusp-native-root-pfaffian} and the root oscillator
proof give complex dimension twenty-four for each \(E_+\) scalar joint
level.  Hermitian adjunction identifies the \(E_-\) level with its complex
conjugate.  If \(v_1,\ldots,v_{24}\) is a complex basis in the \(E_+\)
level, then
\begin{equation*}
 v_j-v_j^\dagger,
 \qquad i(v_j+v_j^\dagger),
 \qquad 1\le j\le24,
\end{equation*}
are forty-eight real linearly independent anti-Hermitian eigenforms and
span the real root eigenspace.  This proves the multiplicity and
\eqref{eq:s6cusp-native-Morse-index}.  Substitution of the exact frequencies
gives \eqref{eq:s6cusp-native-negative-bands-coordinates}; substitution of
their retained cusp expansions gives
\eqref{eq:s6cusp-native-negative-band-cusp}.
\end{proof}

The action excess and the lowest Jacobi branch satisfy the exact identity
\begin{align}
 \nu_{R,0}
 &=-2(\mathcal B_{R,+}-\mathcal B_{R,-})
  =-\frac{4\pi(p_R+6T_R)}{\delta_R},
 \label{eq:s6cusp-native-lowest-instability}\\
 \operatorname{YM}_R(A_{\eta,R})
 -\frac{24\pi^2}{g_{\mathrm{YM}}^2}
 &=\frac{\delta_R}{8g_{\mathrm{YM}}^2}\nu_{R,0}^2.
 \label{eq:s6cusp-native-action-instability}
\end{align}
Indeed, the first equality follows from the exact frequency difference, and
substitution into \eqref{eq:s6cusp-native-action-excess} proves the second.
The spin-one shift and its sign agree with the homogeneous-background
operator \(K_\pm=-D^2\pm2gB\) and its lowest negative branch in
\cite[eqs.~(1)--(3)]{Oxman2010}; the compact bundle, the twenty-four theta
states, the Coulomb equation, and the index \(144\) are proved above rather
than imported from that comparison.

\subsection{One decomposable native operator with continuous threshold zero}
\label{subsec:s6cusp-native-direct-integral}

The marking supplies an exact common Hilbert-space model.  Let
\(\mathbb T_y^4:=\mathbb R_y^4/\mathbb Z^4\), let \(\mathcal L_\eta\) be
the fixed line bundle with clutch factors
\eqref{eq:s6cusp-native-line-clutch-one}--
\eqref{eq:s6cusp-native-line-clutch-two}, and define
\begin{equation}\label{eq:s6cusp-native-marking-diffeomorphism}
 \varphi_R:\mathbb T_y^4\longrightarrow F_R,
 \qquad
 \varphi_R([y]):=[P_Ry].
\end{equation}
Its inverse is \([x]\mapsto[P_R^{-1}x]\), so it is a diffeomorphism and
identifies \(\varphi_R^*\mathcal L_{\eta,R}\) with \(\mathcal L_\eta\).
The pulled-back metric and volume form are
\begin{equation}\label{eq:s6cusp-native-pulled-metric-volume}
 \varphi_R^*g_R=(\dd y)^{\mathsf T}P_R^{\mathsf T}P_R\dd y,
 \qquad
 \varphi_R^*\operatorname{vol}_{g_R}
 =\delta_R\dd y^1\dd y^2\dd y^3\dd y^4.
\end{equation}
Consequently
\begin{equation}\label{eq:s6cusp-native-marking-unitary}
 U_R:L^2(F_R,\mathcal L_{\eta,R};\operatorname{vol}_{g_R})
 \longrightarrow
 \mathscr K:=L^2(\mathbb T_y^4,\mathcal L_\eta;\dd^4y),
 \qquad
 (U_Rs)(y):=\delta_R^{1/2}s(P_Ry)
\end{equation}
is unitary.  Put
\begin{equation}\label{eq:s6cusp-native-fixed-fibre-operator}
 \widehat H_R:=U_RH_{\eta,R}U_R^{-1},
 \qquad
 \operatorname{Dom}(\widehat H_R)
 =H^2(\mathbb T_y^4,\mathcal L_\eta).
\end{equation}
The metric matrix \(P_R^{\mathsf T}P_R\), its inverse, and the fixed
connection coefficients make \(R\mapsto\widehat H_R\) a smooth family from
the common \(H^2\) domain to \(\mathscr K\).

Fix \(R_0\) in the regular radial range and define
\begin{align}
 \mathscr H_\eta
 &:=L^2([R_0,\infty),\dd R;\mathscr K),
 \label{eq:s6cusp-native-direct-integral-Hilbert}\\
 (\mathscr M_\eta\Psi)(R)&:=\widehat H_R\Psi(R),
 \label{eq:s6cusp-native-direct-integral-action}\\
 \operatorname{Dom}(\mathscr M_\eta)
 &:=\left\{\Psi\in\mathscr H_\eta:
 \begin{array}{l}
 \Psi(R)\in H^2(\mathbb T_y^4,\mathcal L_\eta)
 \text{ for almost every }R,\\
 R\mapsto\widehat H_R\Psi(R)\text{ is measurable, and}\\
 \displaystyle\int_{R_0}^{\infty}
 \lVert\widehat H_R\Psi(R)\rVert_{\mathscr K}^2\dd R<\infty
 \end{array}\right\}.
 \label{eq:s6cusp-native-direct-integral-domain}
\end{align}

\begin{theorem}[Self-adjointness and continuous essential threshold]
\label{thm:s6cusp-native-direct-integral-threshold}
The operator \(\mathscr M_\eta\) is self-adjoint and nonnegative, and
\begin{equation}\label{eq:s6cusp-native-direct-integral-spectrum-zero}
 0\in\operatorname{Spec}_{\mathrm{ess}}(\mathscr M_\eta)
 \cap\operatorname{Spec}_{\mathrm c}(\mathscr M_\eta),
 \qquad
 \ker\mathscr M_\eta=\{0\}.
\end{equation}
\end{theorem}

\begin{proof}
The common-domain smoothness in
\eqref{eq:s6cusp-native-fixed-fibre-operator}, the local elliptic estimate,
and the resolvent identity imply local norm continuity of
\(R\mapsto(\widehat H_R\pm i)^{-1}\).  In particular the resolvent field is
strongly measurable and has norm at most one.  Define the bounded operators
on \(\mathscr H_\eta\)
\begin{equation*}
 (\mathscr R_\pm\Psi)(R):=
 (\widehat H_R\pm i)^{-1}\Psi(R).
\end{equation*}
For \(\Phi=\mathscr R_\pm\Psi\),
\begin{equation*}
 \widehat H_R\Phi(R)=\Psi(R)\mp i\Phi(R),
\end{equation*}
so \(\Phi\) belongs to
\(\operatorname{Dom}(\mathscr M_\eta)\) and
\((\mathscr M_\eta\pm i)\Phi=\Psi\).  Fibrewise integration proves
symmetry of \(\mathscr M_\eta\); surjectivity of both
\(\mathscr M_\eta\pm i\) proves self-adjointness.  Equation
\eqref{eq:s6cusp-native-ground-lower-bound} proves nonnegativity.

The rank-six ground projection is isolated at each finite \(R\) by the
strict gap \(2\min\{\mathcal B_{R,+},\mathcal B_{R,-}\}\).  A local Riesz
contour and the norm-continuous resolvent therefore give a smooth rank-six
projection field.  These ranges form a smooth vector bundle over the
interval \([R_0,\infty)\); the interval is contractible, so choose a smooth
unit section \(R\mapsto\phi_R\) with
\begin{equation}\label{eq:s6cusp-native-ground-section}
 \widehat H_R\phi_R=E_{R;0,0}\phi_R,
 \qquad \lVert\phi_R\rVert_{\mathscr K}=1.
\end{equation}
For every integer \(j>R_0\), put
\begin{equation}\label{eq:s6cusp-native-Weyl-sequence}
 \Psi_j(R):=\mathbf1_{[j,j+1]}(R)\phi_R.
\end{equation}
The vectors are orthonormal; hence they converge weakly to zero.  Equations
\eqref{eq:s6cusp-native-ground-section} and
\(E_{R;0,0}=14\pi R^{-1}+O(R^{-2})\) give
\begin{equation}\label{eq:s6cusp-native-Weyl-residual}
 \lVert\mathscr M_\eta\Psi_j\rVert^2
 =\int_j^{j+1}E_{R;0,0}^2\dd R
 =O(j^{-2})\longrightarrow0.
\end{equation}
Thus \((\Psi_j)\) is a singular Weyl sequence at zero, proving
\(0\in\operatorname{Spec}_{\mathrm{ess}}(\mathscr M_\eta)\).

If \(\mathscr M_\eta\Psi=0\), then
\(\widehat H_R\Psi(R)=0\) for almost every \(R\).  The strict fibrewise
bound
\(\widehat H_R\ge E_{R;0,0}I>0\) forces \(\Psi(R)=0\) almost everywhere,
so the kernel is zero.  A self-adjoint operator has no residual spectrum;
zero is therefore in the continuous spectrum, which proves
\eqref{eq:s6cusp-native-direct-integral-spectrum-zero}.
\end{proof}

The exact morphism from the radial family to this operator is the marking
unitary \eqref{eq:s6cusp-native-marking-unitary} followed by the decomposable
action \eqref{eq:s6cusp-native-direct-integral-action}.  There is no
derivative in the \(R\)-direction: a horizontal kinetic completion would
contain an additional differential expression such as \(-\partial_R^2\)
together with a specified measure and boundary domain.  Hence
\cref{thm:s6cusp-native-direct-integral-threshold} proves a continuous
threshold for the displayed decomposable magnetic operator, not the
spectrum of an unconstructed horizontal completion or of a reconstructed
interacting Yang--Mills Hamiltonian.


\subsection{The exceptional elliptic quotient and its exact restriction
from the native class}
\label{subsec:s6cusp-exceptional-elliptic-restriction}

The T1 first-coordinate construction and the native four-dimensional
construction are related by a restriction morphism.  On the complete marked
parameter domain, write
\begin{equation}\label{eq:s6cusp-general-marked-fibre}
 \Pi(z)=
 \begin{pmatrix}
 6\mu(z)&\tau(z)&1&0\\
 \beta(z)&\mu(z)&0&1
 \end{pmatrix},
 \qquad
 F_z:=\C^2/\Pi(z)\Z^4,
\end{equation}
and retain
\begin{align}
 T_z&:=\operatorname{Im}\tau(z)>0,&
 m_z&:=\operatorname{Im}\mu(z),&
 p_z&:=\operatorname{Im}\beta(z),
 \label{eq:s6cusp-general-imaginary-coordinates}\\
 \delta_z&:=6m_z^2-p_zT_z>0,&
 T_z\left(p_z-\frac{6m_z^2}{T_z}\right)&=-\delta_z\ne0.
 \label{eq:s6cusp-general-unchanged-expression}
\end{align}
The corresponding real period matrix is
\begin{equation}\label{eq:s6cusp-general-real-period-matrix}
 P_z=
 \begin{pmatrix}
 6\operatorname{Re}\mu(z)&\operatorname{Re}\tau(z)&1&0\\
 6m_z&T_z&0&0\\
 \operatorname{Re}\beta(z)&\operatorname{Re}\mu(z)&0&1\\
 p_z&m_z&0&0
 \end{pmatrix}.
\end{equation}
Equations \eqref{eq:s6cusp-integral-coframe}--
\eqref{eq:s6cusp-native-line-connection}, with \(P_R\) replaced by
\(P_z\), define \(\eta_z\), \(\mathcal L_{\eta,z}\), and
\(\nabla_{\eta,z}\) on this fibre.

The source uniformisation used below is the holomorphic quotient map
\begin{equation}\label{eq:s6cusp-source-uniformisation}
 \pi:\mathfrak h_z\longrightarrow B\setminus\{p_0\},
 \qquad
 B=\mathbb P^1,
 \qquad
 \mathfrak h_z/\Delta\simeq B\setminus\{p_0\},
\end{equation}
where
\begin{equation}\label{eq:s6cusp-source-deck-presentation}
 \Delta=\langle g_1,g_2\mid g_1^3=g_2^4=1\rangle,
 \qquad
 g_0=(g_1g_2)^{-1}.
\end{equation}
Here \(p_0\) is the distinguished cusp and \(\mathfrak h_z\) is connected;
these are the exact conventions of Proposition~2.11 and
Definition~3.1 of the source transcription \cite{S6Internal}.  In particular,
\(g_0,g_2\) generate \(\Delta\), since
\(g_1=g_0^{-1}g_2^{-1}\).

For \(a,b\in\Z\), define
\begin{equation}\label{eq:s6cusp-descent-loci}
 \mathscr D_{a,b}:=
 \{z\in\mathfrak h_z:6\mu(z)=a\tau(z)+b\},
 \qquad
 \mathscr D:=\bigcup_{(a,b)\in\Z^2}\mathscr D_{a,b}.
\end{equation}

\begin{theorem}[Exact quotient, kernel, and degree-six restriction]
\label{thm:s6cusp-exceptional-elliptic-restriction}
The first-coordinate projection \(\operatorname{pr}_1:\C^2\to\C\)
descends to a surjective holomorphic homomorphism
\begin{equation}\label{eq:s6cusp-first-coordinate-quotient}
 q_{a,b,z}:F_z\longrightarrow
 E_{\tau(z)}:=\C/(\Z\tau(z)+\Z),
 \qquad
 [(\zeta_1,\zeta_2)]\longmapsto[\zeta_1],
\end{equation}
if and only if \(z\in\mathscr D_{a,b}\) for a unique pair
\((a,b)\in\Z^2\).  For such a point, put
\begin{equation}\label{eq:s6cusp-exceptional-kernel-parameter}
 \widehat\gamma_{a,b}:=
 \widehat\gamma-a\widehat u-b\widehat w,
 \qquad
 \tau_{\mathrm f}(z):=a\mu(z)-\beta(z).
\end{equation}
Then
\begin{equation}\label{eq:s6cusp-exceptional-kernel-imaginary-part}
 \operatorname{Im}\tau_{\mathrm f}(z)
 =\frac{\delta_z}{T_z}
 =-\left(
 p_z-\frac{6m_z^2}{T_z}
 \right)>0,
\end{equation}
and there is an exact sequence of complex tori
\begin{equation}\label{eq:s6cusp-exceptional-exact-sequence}
 0\longrightarrow E_{\tau_{\mathrm f}(z)}
 \overset{\iota_{a,b,z}}{\longrightarrow}F_z
 \overset{q_{a,b,z}}{\longrightarrow}E_{\tau(z)}
 \longrightarrow0,
 \qquad
 \iota_{a,b,z}([w])=[(0,w)].
\end{equation}
On integral first homology, in the target basis represented by
\((\tau(z),1)\), the two maps are
\begin{align}
 (q_{a,b,z})_*\widehat\gamma&=(a,b),&
 (q_{a,b,z})_*\widehat u&=(1,0),&
 (q_{a,b,z})_*\widehat w&=(0,1),&
 (q_{a,b,z})_*\widehat\delta&=(0,0),
 \label{eq:s6cusp-exceptional-homology-quotient}\\
 (\iota_{a,b,z})_*[1]&=\widehat\delta,&
 (\iota_{a,b,z})_*[\tau_{\mathrm f}(z)]
 &=-\widehat\gamma_{a,b}.
 \label{eq:s6cusp-exceptional-homology-inclusion}
\end{align}
The native bundle restricts to a degree-six line bundle:
\begin{equation}\label{eq:s6cusp-native-restriction-degree}
 c_1(\iota_{a,b,z}^*\mathcal L_{\eta,z})
 =\iota_{a,b,z}^*\eta_z,
 \qquad
 \int_{E_{\tau_{\mathrm f}(z)}}
 \iota_{a,b,z}^*\eta_z=6.
\end{equation}
\end{theorem}

\begin{proof}
The first coordinates of the four columns of \(\Pi(z)\) are
\begin{equation}\label{eq:s6cusp-first-coordinate-periods}
 6\mu(z),\qquad\tau(z),\qquad1,\qquad0.
\end{equation}
Thus \(\operatorname{pr}_1\) carries the source lattice into
\(\Z\tau(z)+\Z\) exactly when
\(6\mu(z)=a\tau(z)+b\) for integers \(a,b\), proving descent.  The
universal-cover projection is surjective, and the two source periods
\(\tau(z),1\) generate the complete target lattice, so the descended map is
surjective and its kernel is connected.  If two pairs represented the same
point, then
\((a-a')\tau(z)+(b-b')=0\); taking imaginary parts gives \(a=a'\), and
then \(b=b'\).  This proves uniqueness.

The integral basis replacement
\begin{equation}\label{eq:s6cusp-exceptional-unimodular-basis}
 (\widehat\gamma,\widehat u,\widehat w,\widehat\delta)
 \longmapsto
 (\widehat\gamma_{a,b},\widehat u,
  \widehat w,\widehat\delta)
\end{equation}
has determinant one, and
\begin{equation}\label{eq:s6cusp-exceptional-kernel-periods}
 \Pi(z)\widehat\gamma_{a,b}
 =(0,\beta(z)-a\mu(z))=(0,-\tau_{\mathrm f}(z)),
 \qquad
 \Pi(z)\widehat\delta=(0,1).
\end{equation}
More explicitly, the kernel of the lattice homomorphism in
\eqref{eq:s6cusp-exceptional-homology-quotient} consists of the vectors
\begin{align*}
 n_\gamma\widehat\gamma+n_u\widehat u
 +n_w\widehat w+n_\delta\widehat\delta
 \quad\text{with}\quad
 an_\gamma+n_u=0,
 \quad bn_\gamma+n_w=0,
\end{align*}
and hence equals
\(\Z\widehat\gamma_{a,b}\oplus\Z\widehat\delta\).
Equations \eqref{eq:s6cusp-exceptional-kernel-periods} therefore prove the
map, kernel, image, and exactness in
\eqref{eq:s6cusp-exceptional-exact-sequence}, as well as both formulas in
\eqref{eq:s6cusp-exceptional-homology-inclusion}.

Taking imaginary parts of the descent identity gives
\(6m_z=aT_z\).  Consequently
\begin{align*}
 \operatorname{Im}\tau_{\mathrm f}(z)
 &=am_z-p_z
 =\frac{6m_z^2}{T_z}-p_z
 =\frac{\delta_z}{T_z},
\end{align*}
which proves \eqref{eq:s6cusp-exceptional-kernel-imaginary-part} while
retaining \eqref{eq:s6cusp-general-unchanged-expression}.

The positively oriented homology basis of
\(E_{\tau_{\mathrm f}(z)}\) is represented by
\((1,\tau_{\mathrm f}(z))\), because
\(\operatorname{Im}\tau_{\mathrm f}(z)>0\).  Using
\eqref{eq:s6cusp-exceptional-homology-inclusion} and
\(\eta=u\wedge w+6\gamma\wedge\delta\), one obtains
\begin{align*}
 \int_{E_{\tau_{\mathrm f}(z)}}\iota_{a,b,z}^*\eta_z
 &=\eta_z(\widehat\delta,-\widehat\gamma_{a,b})\\
 &=\eta_z\bigl(
 \widehat\delta,
 -\widehat\gamma+a\widehat u+b\widehat w
 \bigr)=6.
\end{align*}
Naturality of the first Chern class under pullback gives the first equality
in \eqref{eq:s6cusp-native-restriction-degree}, completing the proof.
\end{proof}

The restriction relation also recovers the T1 Landau operator with all
coefficients fixed.  Write
\(w=x^3+ix^4\) on the kernel line in the physical covering coordinates.
Restricting \eqref{eq:s6cusp-native-physical-matrix} to
\(x^1=x^2=0\) and using
\(\delta_z=T_z\operatorname{Im}\tau_{\mathrm f}(z)\) gives
\begin{equation}\label{eq:s6cusp-exceptional-restricted-curvature}
 \iota_{a,b,z}^*\eta_z
 =\frac{6T_z}{\delta_z}\,\dd x^3\wedge\dd x^4
 =\frac{6}{\operatorname{Im}\tau_{\mathrm f}(z)}
 \dd x^3\wedge\dd x^4.
\end{equation}
Thus the real magnetic field of the restricted unitary connection is
\begin{equation}\label{eq:s6cusp-exceptional-magnetic-field}
 \mathcal B_{\mathrm f}(z)
 =\frac{12\pi}{\operatorname{Im}\tau_{\mathrm f}(z)}.
\end{equation}

\begin{corollary}[The T1 degree-six spectrum as an exact restriction]
\label{cor:s6cusp-exceptional-landau-restriction}
Let
\begin{equation}\label{eq:s6cusp-exceptional-landau-Hamiltonian}
 H_{\mathrm f,z}:=\frac12
 (\iota_{a,b,z}^*\nabla_{\eta,z})^*
 (\iota_{a,b,z}^*\nabla_{\eta,z})
\end{equation}
be the Friedrichs operator on
\(L^2(E_{\tau_{\mathrm f}(z)},
\iota_{a,b,z}^*\mathcal L_{\eta,z})\), with the metric induced from the
vertical complex line in \(\C^2\).  Then
\begin{equation}\label{eq:s6cusp-exceptional-landau-spectrum}
 \operatorname{Spec}(H_{\mathrm f,z})
 =\left\{
 \frac{12\pi}{\operatorname{Im}\tau_{\mathrm f}(z)}
 \left(n+\frac12\right):n\in\Z_{\ge0}
 \right\}.
\end{equation}
Every displayed level has complex multiplicity six.  Subtracting the exact
ground energy
\(6\pi/\operatorname{Im}\tau_{\mathrm f}(z)\) produces levels
\begin{equation}\label{eq:s6cusp-exceptional-ground-shifted-spectrum}
 \left\{
 \frac{12\pi n}{\operatorname{Im}\tau_{\mathrm f}(z)}:
 n\in\Z_{\ge0}
 \right\},
\end{equation}
with a six-dimensional kernel and adjacent-level spacing
\(12\pi/\operatorname{Im}\tau_{\mathrm f}(z)\).
\end{corollary}

\begin{proof}
Let \(\Pi_3,\Pi_4\) be the self-adjoint covariant momenta on the restricted
bundle.  Equations \eqref{eq:s6cusp-exceptional-restricted-curvature}--
\eqref{eq:s6cusp-exceptional-magnetic-field} give
\begin{equation}\label{eq:s6cusp-exceptional-momentum-commutator}
 [\Pi_3,\Pi_4]=i\mathcal B_{\mathrm f}(z).
\end{equation}
Define
\begin{equation}\label{eq:s6cusp-exceptional-ladder}
 \mathfrak a_{\mathrm f}
 :=\frac{\Pi_3+i\Pi_4}{\sqrt{2\mathcal B_{\mathrm f}(z)}},
 \qquad
 \mathfrak a_{\mathrm f}^*
 :=\frac{\Pi_3-i\Pi_4}{\sqrt{2\mathcal B_{\mathrm f}(z)}}.
\end{equation}
Direct multiplication gives
\begin{equation}\label{eq:s6cusp-exceptional-oscillator-factorization}
 [\mathfrak a_{\mathrm f},\mathfrak a_{\mathrm f}^*]=I,
 \qquad
 H_{\mathrm f,z}
 =\mathcal B_{\mathrm f}(z)
 \left(\mathfrak a_{\mathrm f}^*\mathfrak a_{\mathrm f}
 +\frac12I\right).
\end{equation}
The lowering argument used for the two native number operators above proves
that the number spectrum is contained in \(\Z_{\ge0}\), and the normalized
raising and lowering maps identify every level with the ground space.  The
ground equation \(\mathfrak a_{\mathrm f}s=0\) is the covariant
\(\bar\partial\)-equation for the standard complex structure of
\(E_{\tau_{\mathrm f}(z)}\).  By
\eqref{eq:s6cusp-native-restriction-degree}, the holomorphic bundle has
degree six.  Riemann--Roch on the elliptic curve gives
\(h^0-h^1=6\), and Serre duality identifies \(h^1\) with the space of
sections of the degree \(-6\) dual bundle, which is zero: a nonzero section
would have an effective zero divisor of negative degree.  Hence the ground
space has dimension six.  Substitution of
\eqref{eq:s6cusp-exceptional-magnetic-field} into
\eqref{eq:s6cusp-exceptional-oscillator-factorization} proves every level,
multiplicity, shift, kernel, and spacing asserted.
\end{proof}

\begin{theorem}[Exact transformation and finiteness of the descent locus]
\label{thm:s6cusp-descent-locus-finite}
The set \(\mathscr D\) is invariant under the source deck group.  More
precisely, the source laws for \(g_0\) and \(g_2\) induce the bijections
\begin{align}
 g_0:\mathscr D_{a,b}&\longrightarrow\mathscr D_{a,a+b},
 & (a,b)&\longmapsto(a,a+b),
 \label{eq:s6cusp-descent-pair-g0}\\
 g_2:\mathscr D_{a,b}&\longrightarrow\mathscr D_{-b,a+6},
 & (a,b)&\longmapsto(-b,a+6).
 \label{eq:s6cusp-descent-pair-g2}
\end{align}
There is a punctured neighbourhood of the distinguished cusp whose inverse
image contains no point of \(\mathscr D\), and the image
\begin{equation}\label{eq:s6cusp-descent-finite-image}
 \pi(\mathscr D)\subset B\setminus\{p_0\}
\end{equation}
is finite.
\end{theorem}

\begin{proof}
The directly sourced transformation laws are
\begin{equation}\label{eq:s6cusp-descent-source-laws}
 \tau\circ g_0=\tau-1,
 \quad \mu\circ g_0=\mu,
 \qquad
 \tau\circ g_2=-\frac1\tau,
 \quad \mu\circ g_2=1+\frac\mu\tau.
\end{equation}
If \(6\mu=a\tau+b\), then
\begin{align*}
 6\mu(g_0z)
 &=a\tau(g_0z)+(a+b),\\
 6\mu(g_2z)
 &=6+a+\frac b\tau
 =(-b)\tau(g_2z)+(a+6).
\end{align*}
The affine maps on pairs have integral inverses
\begin{equation}\label{eq:s6cusp-descent-pair-inverses}
 (a,b)\longmapsto(a,b-a),
 \qquad
 (a,b)\longmapsto(b-6,-a),
\end{equation}
respectively.  Since \(g_0,g_2\) generate the source deck group, these
formulas prove invariance and the asserted bijections.

No function
\begin{equation}\label{eq:s6cusp-descent-holomorphic-function}
 f_{a,b}:=6\mu-a\tau-b
\end{equation}
vanishes identically.  Indeed, applying \(g_0\) to an identity
\(f_{a,b}=0\) and subtracting the original identity gives \(a=0\).
Then \(6\mu=b\) is constant, whereas applying \(g_2\) gives
\begin{equation}\label{eq:s6cusp-descent-constant-contradiction}
 b=6+\frac b\tau,
 \qquad
 (b-6)\tau=b,
\end{equation}
which contradicts the nonconstancy of \(\tau\) proved by its source modular
relation.

Suppose points of \(\mathscr D\) occurred arbitrarily close to the cusp in
the quotient.  Invariance permits representatives
\(z_j\in\mathscr D_{a_j,b_j}\) in the distinguished cusp component with
\(t_c(z_j)\to0\).  Taking imaginary parts gives the exact integer identity
\begin{equation}\label{eq:s6cusp-descent-integer-a}
 a_j=\frac{6\operatorname{Im}\mu(z_j)}
 {\operatorname{Im}\tau(z_j)}.
\end{equation}
The numerator remains bounded and the denominator tends to infinity, so
\(a_j=0\) for all sufficiently large \(j\).  Then
\(6\mu(z_j)=b_j\in\Z\).  Boundedness of \(\mu\) makes the integer sequence
\((b_j)\) bounded; after taking a subsequence it is one constant \(b\).
The source function \(\mu\) descends on the cusp component to a holomorphic
function of \(t_c\) extending through \(t_c=0\).  Its function
 \(6\mu(t_c)-b\) has zeros accumulating at that interior point, so it
 vanishes identically on the cusp component.  Pulling the identity back to
 the connected open subset of \(\mathfrak h_z\) represented by that component
 and applying analytic continuation on the connected surface
 \(\mathfrak h_z\), one obtains \(f_{0,b}\equiv0\) on all of
 \(\mathfrak h_z\).  This contradicts
 \eqref{eq:s6cusp-descent-holomorphic-function} and proves the cusp exclusion.

Delete a smaller cusp horodisc from a closed source fundamental triangle.
The remainder \(K\subset\mathfrak h_z\) is compact and represents every
point of the quotient outside the excluded cusp neighbourhood.  Define the
finite constants
\begin{equation}\label{eq:s6cusp-descent-compact-bounds}
 M_\mu:=\max_K|\mu|,
 \qquad
 M_\tau:=\max_K|\tau|,
 \qquad
 \varepsilon_\tau:=\min_K\operatorname{Im}\tau>0.
\end{equation}
At a zero of \(f_{a,b}\) in \(K\), equations
\begin{equation}\label{eq:s6cusp-descent-pair-bounds}
 |a|\le\frac{6M_\mu}{\varepsilon_\tau},
 \qquad
 |b|=|6\mu-a\tau|
 \le6M_\mu+|a|M_\tau
\end{equation}
leave only finitely many integer pairs.  For each retained pair,
\(f_{a,b}\) is a nonzero holomorphic function on a neighbourhood of the
compact set \(K\), so its zeros in \(K\) are finite by the identity theorem.
The union over the finite pair set is finite, and the cusp exclusion proves
\eqref{eq:s6cusp-descent-finite-image}.
\end{proof}

The exact conclusion is therefore twofold.  The native four-dimensional
bundle \(\mathcal L_{\eta,z}\) and its spectrum do not require membership in
\(\mathscr D\).  Membership in \(\mathscr D\) is exactly what supplies the
additional quotient map \eqref{eq:s6cusp-first-coordinate-quotient}; on that
finite locus, the exceptional T1 degree-six bundle and Landau operator are
the restriction in \eqref{eq:s6cusp-native-restriction-degree}--
\eqref{eq:s6cusp-exceptional-landau-Hamiltonian}.  Since
\(\mathscr D\) is absent from a cusp neighbourhood, the exceptional
two-dimensional spacing in
\eqref{eq:s6cusp-exceptional-ground-shifted-spectrum} has no cusp sequence,
whereas the native four-dimensional gap
\eqref{eq:s6cusp-native-gap-exact-eventual} does tend to zero.

\subsection{The native \texorpdfstring{\(T^2\times\R^2\)}{T2 x R2}
pointed limit}

The fixed third and fourth columns of \(P_R\) make the native cusp
two-dimensional at infinity unless those two cycles are also unrolled.
This subsection constructs the unrolling as an actual finite holomorphic
cover and proves its compatibility with the source monodromy.

\begin{proposition}[Universal cover, quotient torus, and native pointed limit]
\label{prop:s6cusp-native-limit}
The quotient map
\begin{equation}\label{eq:s6cusp-universal-cover}
 \upsilon_R:\R^4\longrightarrow F_R,\qquad x\longmapsto[x]
\end{equation}
is the universal covering map and
\[
 \ker\upsilon_R=P_R\Z^4.
\]
Under \(\R^4=\C^2\), it is holomorphic.  Define
\begin{equation}\label{eq:s6cusp-BR-matrix}
 B_R:=
 \begin{pmatrix}
 6m_R&T_R\\
 p_R&m_R
 \end{pmatrix},
 \qquad
 \Sigma_R:=\R^2/B_R\Z^2,
\end{equation}
and
\begin{equation}\label{eq:s6cusp-fixed-subtorus}
 K_R:=
 \left\{[(x^1,0,x^3,0)]:
 (x^1,x^3)\in\R^2\right\}\subset F_R.
\end{equation}
Then
\begin{equation}\label{eq:s6cusp-native-exact-sequence}
 0\longrightarrow K_R
 \longrightarrow F_R
 \overset{\pi_R}{\longrightarrow}\Sigma_R
 \longrightarrow0,
 \qquad
 \pi_R([x])=[(x^2,x^4)]
\end{equation}
is an exact sequence of real tori, with
\[
 K_R\cong\R^2/\Z^2.
\]
In particular,
\begin{equation}\label{eq:s6cusp-native-systole-upper}
 \operatorname{sys}(F_R)\le1,\qquad
 \operatorname{inj}(F_R)\le\frac12.
\end{equation}

Let
\[
 F_\infty^{\mathrm{native}}
 :=
 \R^4/
 \Z\langle(1,0,0,0),(0,0,1,0)\rangle
 \cong T^2\times\R^2
\]
with its quotient Euclidean metric and basepoint zero.  Then
\begin{equation}\label{eq:s6cusp-native-pointed-limit}
 (F_R,[0])\longrightarrow
 (F_\infty^{\mathrm{native}},[0])
\end{equation}
in the pointed smooth and pointed local-isometric senses as
\(R\to\infty\).
\end{proposition}

\begin{proof}
The first assertion is the defining quotient of
\eqref{eq:s6cusp-flat-fibre}; the action of \(P_R\Z^4\) on \(\R^4\) is
free and properly discontinuous, and \(\R^4\) is simply connected.
The same quotient is \(\C^2/\Pi_R\Lambda\), so the quotient map is
holomorphic.

Projection to the second and fourth real coordinates sends the first two
columns of \(P_R\) to the two columns of \(B_R\), and sends its third and
fourth columns to zero.  Hence \(\pi_R\) is well defined and surjective.
If \(\pi_R([x])=0\), there is \(n\in\Z^2\) for which
\((x^2,x^4)^{\mathsf T}=B_Rn\).  Subtracting from \(x\) the corresponding
integer combination of the first two columns of \(P_R\) leaves
\((v^1,0,v^3,0)\), proving \(\ker\pi_R=K_R\).
If \((v^1,0,v^3,0)\in P_R\Z^4\), invertibility of \(B_R\), whose
determinant is \(\delta_R>0\), forces the first two lattice coefficients
to vanish.  The remaining two columns are
\((1,0,0,0)\) and \((0,0,1,0)\), so the kernel of
\(\R^2\to K_R\) is exactly \(\Z^2\).  This proves every map, kernel, and
image in \eqref{eq:s6cusp-native-exact-sequence}.  The two displayed unit
periods prove \eqref{eq:s6cusp-native-systole-upper}; for a flat torus the
injectivity radius is one half its shortest period.

Reorder the covering coordinates as
\((x^1,x^3;x^2,x^4)\), and put
\begin{equation}\label{eq:s6cusp-AR-matrix}
 A_R:=
 \begin{pmatrix}
 6a_R&u_R\\
 c_R^{\mathrm{re}}&a_R
 \end{pmatrix}.
\end{equation}
Every native period has the exact form
\begin{equation}\label{eq:s6cusp-native-period-block-form}
 \lambda_R(n,m)=
 \begin{pmatrix}A_Rn+m\\B_Rn\end{pmatrix},
 \qquad n,m\in\Z^2.
\end{equation}
By \eqref{eq:s6cusp-imaginary-identities},
\begin{equation}\label{eq:s6cusp-BR-cusp-splitting}
 B_R=RJ+C_R,\qquad
 J=\begin{pmatrix}0&1\\-1&0\end{pmatrix},\qquad
 C_R=
 \begin{pmatrix}
 6m_R&h_R\\
 d_R-h_R&m_R
 \end{pmatrix}.
\end{equation}
Holomorphy at \(q_R=0\) makes \(A_R\) and \(C_R\) bounded.  Choose
\begin{equation}\label{eq:s6cusp-Cstar}
 C_*:=\sup_{R\ge R_0}\lVert C_R\rVert_{\mathrm{op}}<\infty.
\end{equation}
Since \(J\) is orthogonal,
\begin{equation}\label{eq:s6cusp-BR-lower-bound}
 \lVert B_Rn\rVert
 \ge(R-C_*)\lVert n\rVert.
\end{equation}
Thus every period not in the fixed lattice
\(\Z\langle(1,0,0,0),(0,0,1,0)\rangle\) has length at least
\(R-C_*\).  On each fixed-radius ball, all such translations disappear
for sufficiently large \(R\), while the two fixed translations remain
exactly.  Quotient distances on that ball therefore agree with the
distances in \(F_\infty^{\mathrm{native}}\), which proves
\eqref{eq:s6cusp-native-pointed-limit}.
\end{proof}

\subsection{A decomposable \texorpdfstring{\(SU(2)\)}{SU(2)}
Yang--Mills saddle and its exact negative mode}
\label{subsec:s6cusp-reducible-su2-saddle}

The third-turn calculation uses a second integral class on the same marked
four-torus.  Its relationship to the native class \(\eta_R\) will be computed
at the end of this subsection.  Equations
\eqref{eq:s6cusp-integral-coframe} and
\eqref{eq:s6cusp-soft-dual-block} give
\begin{align}
 \vartheta_R^1
 &=\frac{m_R\,\dd x^2-T_R\,\dd x^4}{\delta_R},&
 \vartheta_R^2
 &=\frac{-p_R\,\dd x^2+6m_R\,\dd x^4}{\delta_R},
 \label{eq:s6cusp-decomposable-coframe}\\
 \xi_R:=\vartheta_R^1\wedge\vartheta_R^2
 &=\frac1{\delta_R}\,\dd x^2\wedge\dd x^4.
 \label{eq:s6cusp-decomposable-flux}
\end{align}
The coefficient in \eqref{eq:s6cusp-decomposable-flux} is the exact
determinant
\[
 \frac{6m_R^2-p_RT_R}{\delta_R^2}=\frac1{\delta_R};
\]
no period entry has been removed.

\begin{proposition}[Explicit line bundle, null clutch, and trivial
\(SU(2)\) carrier]
\label{prop:s6cusp-su2-null-clutch}
On \(\R^4_y\times\C\), impose
\begin{align}
 (y,z)&\sim(y+e_1,e^{-2\pi i y^2}z),
 \label{eq:s6cusp-decomposable-line-clutch}\\
 (y,z)&\sim(y+e_j,z),\qquad j=2,3,4.
 \label{eq:s6cusp-decomposable-line-trivial-clutches}
\end{align}
The quotient is a line bundle \(\mathcal M_R\to F_R\), and
\begin{equation}\label{eq:s6cusp-decomposable-line-connection}
 \nabla^{\mathcal M_R}:=\dd+2\pi i y^1\dd y^2
\end{equation}
descends with
\begin{equation}\label{eq:s6cusp-decomposable-line-curvature}
 F_{\nabla^{\mathcal M_R}}=2\pi i\xi_R,
 \qquad
 c_1(\mathcal M_R)=-[\xi_R]
\end{equation}
under the convention \(c_1=[iF/(2\pi)]\).

The determinant-one bundle
\begin{equation}\label{eq:s6cusp-reducible-su2-bundle}
 \mathcal E_R:=\mathcal M_R\oplus\mathcal M_R^{-1}
\end{equation}
is a topologically trivial \(SU(2)\) bundle.  More precisely, its only
nonidentity clutch map is
\begin{equation}\label{eq:s6cusp-su2-clutch-loop}
 G(v):=
 \begin{pmatrix}e^{-2\pi i v}&0\\0&e^{2\pi i v}\end{pmatrix},
 \qquad v\in\R/\Z,
\end{equation}
and, under \(SU(2)\simeq S^3\), the following displayed homotopy contracts
it to the identity:
\begin{equation}\label{eq:s6cusp-quaternion-su2-map}
 \rho(a+b\mathbf i+c\mathbf j+d\mathbf k)
 :=
 \begin{pmatrix}
  a+di&c+bi\\
  -c+bi&a-di
 \end{pmatrix}.
\end{equation}
For a unit quaternion this matrix is in \(SU(2)\), and
\(\rho(\cos(2\pi v)-\mathbf k\sin(2\pi v))=G(v)\).  Define the
quaternion-valued map
\begin{equation}\label{eq:s6cusp-su2-explicit-null-homotopy-quaternion}
 \mathfrak q(s,v):=
 \begin{cases}
  \cos(\pi s)
  \bigl(\cos(2\pi v)-\mathbf k\sin(2\pi v)\bigr)
  +\sin(\pi s)\mathbf i,
  &0\le s\le\frac12,\\[1mm]
  \cos\!\left(\dfrac{\pi(2s-1)}2\right)\mathbf i
  +\sin\!\left(\dfrac{\pi(2s-1)}2\right)1,
  &\frac12\le s\le1.
 \end{cases}
\end{equation}
The required matrix-valued clutch homotopy is the typed composite
\begin{equation}\label{eq:s6cusp-su2-explicit-null-homotopy}
 \mathscr K:[0,1]\times(\R/\Z)\longrightarrow SU(2),
 \qquad
 \mathscr K(s,v):=\rho\bigl(\mathfrak q(s,v)\bigr).
\end{equation}
It obeys \(\mathscr K(0,v)=G(v)\),
\(\mathscr K(1,v)=I_2\), and is periodic in \(v\).
\end{proposition}

\begin{proof}
The clutch factors commute and are unchanged by the three translations for
which they are declared trivial, so they define the quotient line bundle.
For a coefficient section one has
\(s(y+e_1)=e^{-2\pi i y^2}s(y)\).  The connection one-form obeys
\[
 2\pi i(y^1+1)\dd y^2
 =2\pi i y^1\dd y^2
 -(e^{-2\pi i y^2})^{-1}\dd(e^{-2\pi i y^2}),
\]
which is exactly its descent law.  Differentiation gives
\(F=2\pi i\dd y^1\wedge\dd y^2=2\pi i\xi_R\), and the stated Chern
class follows from the fixed convention.

The determinant of \(G(v)\) is one.  Both pieces of
\eqref{eq:s6cusp-su2-explicit-null-homotopy-quaternion} have quaternion norm one; at
\(s=1/2\) they both equal \(\mathbf i\), so the formula is continuous.
The first piece starts at the quaternion corresponding to
\eqref{eq:s6cusp-su2-clutch-loop}, and the second ends at one.
For \(a^2+b^2+c^2+d^2=1\), direct multiplication gives
\(\rho(q)^*\rho(q)=I\) and
\(\det\rho(q)=a^2+b^2+c^2+d^2=1\), which verifies the stated
quaternion--matrix identification on every point of the homotopy.

For an explicit bundle trivialisation, use the base marking
\begin{equation}\label{eq:s6cusp-Sigma-marking}
 \overline\Phi_R:\R^2/\Z^2\longrightarrow\Sigma_R,
 \qquad [u,v]\longmapsto[B_R(u,v)^{\mathsf T}],
\end{equation}
whose inverse is induced by \(B_R^{-1}\).  The bundle
\(\mathcal E_R\) is the pullback under \(\pi_R\) of the rank-two bundle
on \(\Sigma_R\) obtained from the cylinder
\([0,1]\times S^1\times\C^2\) by
\((0,v,z)\sim(1,v,G(v)z)\).  The map
\[
 [s,v,z]\longmapsto
 \bigl([s,v],\mathscr K(s,v)z\bigr)
\]
is well defined at the glued boundary because
\(\mathscr K(0,v)=G(v)\) and \(\mathscr K(1,v)=1\); its fibrewise inverse
uses \(\mathscr K(s,v)^{-1}\).  It is therefore an explicit trivialisation
of the base bundle, and pulling it back by \(\pi_R\) trivialises
\(\mathcal E_R\).  In agreement with this construction,
\[
 c_2(\mathcal E_R)
 =-c_1(\mathcal M_R)^2=-[\xi_R]^2=0
\]
because \(\xi_R\wedge\xi_R=0\).
\end{proof}

Put
\begin{equation}\label{eq:s6cusp-su2-Cartan-generator}
 H:=\begin{pmatrix}i&0\\0&-i\end{pmatrix},
 \qquad
 \langle X,Y\rangle_{\mathfrak{su}(2)}
 :=-\frac12\operatorname{Tr}(XY),
\end{equation}
so that \(\langle H,H\rangle_{\mathfrak{su}(2)}=1\).  In the diagonal
clutch frame, the direct-sum connection on \(\mathcal E_R\) has
connection one-form and curvature
\begin{align}
 A_R^{\mathrm{red}}&:=2\pi y^1\dd y^2\,H,
 \label{eq:s6cusp-reducible-connection}\\
 F_R^{\mathrm{red}}
 &=2\pi\xi_R H
 =\frac{2\pi}{\delta_R}\dd x^2\wedge\dd x^4\,H.
 \label{eq:s6cusp-reducible-curvature}
\end{align}

\begin{theorem}[Exact nonflat Yang--Mills critical point and action]
\label{thm:s6cusp-reducible-yang-mills-critical}
The connection \(A_R^{\mathrm{red}}\) is nonflat and satisfies
\begin{equation}\label{eq:s6cusp-reducible-yang-mills-equation}
 \dd_{A_R^{\mathrm{red}}}F_R^{\mathrm{red}}=0,
 \qquad
 \dd_{A_R^{\mathrm{red}}}^*F_R^{\mathrm{red}}=0.
\end{equation}
Its Yang--Mills action in the convention
\eqref{eq:s6cusp-yang-mills-action} is
\begin{align}
 \operatorname{YM}_R(A_R^{\mathrm{red}})
 &=\frac{2\pi^2}{g_{\mathrm{YM}}^2\delta_R}
 \label{eq:s6cusp-reducible-action-delta}\\
 &=-\frac{2\pi^2}{g_{\mathrm{YM}}^2
 \operatorname{Im}\tau_R
 \left(
  \operatorname{Im}\beta_R-
  \dfrac{6(\operatorname{Im}\mu_R)^2}
  {\operatorname{Im}\tau_R}
 \right)}.
 \label{eq:s6cusp-reducible-action-unchanged}
\end{align}
Along the source radial cusp it has the two-term expansion
\begin{equation}\label{eq:s6cusp-reducible-action-cusp}
 \operatorname{YM}_R(A_R^{\mathrm{red}})
 =\frac{2\pi^2}{g_{\mathrm{YM}}^2R^2}
 \left[
  1-\frac{2h_0-d_0}{R}+O(R^{-2})
 \right].
\end{equation}
\end{theorem}

\begin{proof}
The coefficient in \eqref{eq:s6cusp-reducible-curvature} is nonzero because
\(\delta_R>0\).  It is constant in the physical covering coordinates at
fixed \(R\), and both \(A_R^{\mathrm{red}}\) and
\(F_R^{\mathrm{red}}\) take values in the one-dimensional Cartan algebra
\(\R H\).  Hence the exterior derivatives of the curvature and its Hodge
dual vanish and all commutator terms vanish.  This proves both equations in
\eqref{eq:s6cusp-reducible-yang-mills-equation}.

The Euclidean norm of \(\dd x^2\wedge\dd x^4\,H\) is one and
\(\operatorname{Vol}(F_R,g_R)=|\det P_R|=\delta_R\).  Direct integration
therefore gives
\[
 \frac1{2g_{\mathrm{YM}}^2}
 \left(\frac{2\pi}{\delta_R}\right)^2\delta_R
 =\frac{2\pi^2}{g_{\mathrm{YM}}^2\delta_R},
\]
which is \eqref{eq:s6cusp-reducible-action-delta}.  Equation
\eqref{eq:s6cusp-unchanged-expression-delta} proves
\eqref{eq:s6cusp-reducible-action-unchanged} without changing the
parenthesis.  Finally, invert the complete expansion
\eqref{eq:s6cusp-delta-two-term-expansion}; the coefficient of \(R^{-1}\)
in the bracket is exactly \(-(2h_0-d_0)\), proving
\eqref{eq:s6cusp-reducible-action-cusp}.
\end{proof}

To construct the negative mode rather than infer it from the shape of a
known formula, define the two physical base periods and their modulus by
\begin{equation}\label{eq:s6cusp-Sigma-complex-periods}
 \omega_{1,R}:=6m_R+ip_R,
 \qquad
 \omega_{2,R}:=T_R+im_R,
 \qquad
 \tau_{\Sigma,R}:=\frac{\omega_{2,R}}{\omega_{1,R}}.
\end{equation}
Since \(\omega_{1,R}\ne0\), direct multiplication gives
\begin{equation}\label{eq:s6cusp-Sigma-modulus-positive}
 \operatorname{Im}\tau_{\Sigma,R}
 =\frac{6m_R^2-p_RT_R}{|\omega_{1,R}|^2}
 =\frac{\delta_R}{|\omega_{1,R}|^2}>0.
\end{equation}
In the base marking coordinates \((u,v)=(y^1,y^2)\),
\begin{equation}\label{eq:s6cusp-Sigma-complex-coordinate}
 z_{\Sigma}:=x^2+ix^4
 =\omega_{1,R}(u+\tau_{\Sigma,R}v).
\end{equation}
The \(E_+\)-root connection induced by
\eqref{eq:s6cusp-reducible-connection} is the connection on
\(\mathcal M_R^2\) with
\begin{equation}\label{eq:s6cusp-root-covariant-derivatives}
 D_u^+:=\partial_u,
 \qquad
 D_v^+:=\partial_v+4\pi i u.
\end{equation}

\begin{proposition}[Two explicit lowest root-Landau sections]
\label{prop:s6cusp-explicit-root-ground-sections}
For \(r\in\{0,1\}\), the normally convergent series
\begin{equation}\label{eq:s6cusp-root-ground-theta-series}
 \phi_{r,R}(u,v):=
 \sum_{k\in\Z}
 \exp\!\left(
  \frac{2\pi i}{\overline{\tau_{\Sigma,R}}}
  \bigl((u+k)^2+r(u+k)\bigr)
  +2\pi i(r+2k)v
 \right)
\end{equation}
obey
\begin{align}
 \phi_{r,R}(u+1,v)&=e^{-4\pi i v}\phi_{r,R}(u,v),&
 \phi_{r,R}(u,v+1)&=\phi_{r,R}(u,v),
 \label{eq:s6cusp-root-ground-automorphy}\\
 (D_v^+-\overline{\tau_{\Sigma,R}}D_u^+)
 \phi_{r,R}&=0.
 \label{eq:s6cusp-root-ground-equation}
\end{align}
They form a basis of the complete solution space of
\eqref{eq:s6cusp-root-ground-equation} with automorphy
\eqref{eq:s6cusp-root-ground-automorphy}.

Let \(D_2^+,D_4^+\) be the same connection in the physical coordinates
\((x^2,x^4)\), and put
\begin{equation}\label{eq:s6cusp-root-magnetic-field}
 \mathcal B_R^{\mathrm{root}}:=\frac{4\pi}{\delta_R}.
\end{equation}
Then each pullback \(\pi_R^*\phi_{r,R}\) satisfies
\begin{align}
 (D_2^+-iD_4^+)\pi_R^*\phi_{r,R}&=0,
 \label{eq:s6cusp-root-physical-ground-equation}\\
 -\sum_{\mu=1}^4(D_\mu^+)^2\pi_R^*\phi_{r,R}
 &=\mathcal B_R^{\mathrm{root}}\pi_R^*\phi_{r,R}.
 \label{eq:s6cusp-root-scalar-ground-energy}
\end{align}
\end{proposition}

\begin{proof}
The real part of the quadratic coefficient in the exponent is
\[
 \operatorname{Re}\frac{2\pi i}{\overline{\tau_{\Sigma,R}}}
 =-\frac{2\pi\operatorname{Im}\tau_{\Sigma,R}}
 {|\tau_{\Sigma,R}|^2}<0.
\]
The series and every termwise derivative consequently converge uniformly on
compact subsets of \(\R^2\).  Replacing \(k\) by \(k+1\) proves the first
automorphy law, while every frequency \(r+2k\) is integral, proving the
second.  On each summand, \(D_v^+\) multiplies by
\(2\pi i(r+2k+2u)\), and
\(\overline{\tau_{\Sigma,R}}D_u^+\) multiplies by that same quantity.
This proves \eqref{eq:s6cusp-root-ground-equation} term by term.

Conversely, expand any smooth solution as
\(\phi(u,v)=\sum_{n\in\Z}f_n(u)e^{2\pi inv}\).  The equation gives the
complete scalar ordinary differential equation
\begin{equation}\label{eq:s6cusp-root-fourier-ode}
 f_n'(u)=\frac{2\pi i(n+2u)}
 {\overline{\tau_{\Sigma,R}}}f_n(u),
\end{equation}
and hence
\[
 f_n(u)=C_n
 \exp\!\left(
  \frac{2\pi i}{\overline{\tau_{\Sigma,R}}}(nu+u^2)
 \right).
\]
The first automorphy law is equivalent to
\begin{equation}\label{eq:s6cusp-root-coefficient-recurrence}
 C_{n+2}=C_n
 \exp\!\left(
  \frac{2\pi i(n+1)}{\overline{\tau_{\Sigma,R}}}
 \right).
\end{equation}
Thus all coefficients are determined by \((C_0,C_1)\).  Choosing
\((C_0,C_1)=(1,0)\) and \((0,1)\) yields precisely the two series in
\eqref{eq:s6cusp-root-ground-theta-series}.  Their Fourier supports have
opposite parity, so they are nonzero and linearly independent; they are
therefore a basis.

Equation \eqref{eq:s6cusp-Sigma-complex-coordinate} gives
\begin{equation}\label{eq:s6cusp-root-coordinate-derivative-map}
 D_2^+-iD_4^+
 =\frac{2}{\omega_{1,R}
 (\tau_{\Sigma,R}-\overline{\tau_{\Sigma,R}})}
 (D_v^+-\overline{\tau_{\Sigma,R}}D_u^+),
\end{equation}
which proves \eqref{eq:s6cusp-root-physical-ground-equation}.  The root
curvature is
\[
 [D_2^+,D_4^+]
 =\frac{4\pi i}{\delta_R}
 =i\mathcal B_R^{\mathrm{root}}.
\]
With \(\Pi_j^+:=-iD_j^+\), define
\[
 \mathfrak a_R^+
 :=\frac{\Pi_2^+-i\Pi_4^+}
 {\sqrt{2\mathcal B_R^{\mathrm{root}}}},
 \qquad
 (\mathfrak a_R^+)^*
 :=\frac{\Pi_2^++i\Pi_4^+}
 {\sqrt{2\mathcal B_R^{\mathrm{root}}}}.
\]
Direct multiplication gives
\begin{equation}\label{eq:s6cusp-root-scalar-factorisation}
 [\mathfrak a_R^+,(\mathfrak a_R^+)^*]=I,
 \qquad
 -(D_2^+)^2-(D_4^+)^2
 =2\mathcal B_R^{\mathrm{root}}
 \left((\mathfrak a_R^+)^*\mathfrak a_R^++\frac12I\right).
\end{equation}
The pullbacks are constant and covariantly constant in the
\(x^1,x^3\) directions.  Their lowering equation and
\eqref{eq:s6cusp-root-scalar-factorisation} prove
\eqref{eq:s6cusp-root-scalar-ground-energy}.
\end{proof}

Let
\begin{equation}\label{eq:s6cusp-root-matrices}
 E_+:=\begin{pmatrix}0&1\\0&0\end{pmatrix},
 \qquad
 E_-:=\begin{pmatrix}0&0\\1&0\end{pmatrix}.
\end{equation}
Then
\begin{equation}\label{eq:s6cusp-adjoint-root-splitting}
 \operatorname{ad}(\mathcal E_R)_{\C}
 =\underline\C H\oplus\mathcal M_R^2E_+
 \oplus\mathcal M_R^{-2}E_-,
 \qquad
 [H,E_+]=2iE_+,
 \quad [H,E_-]=-2iE_-.
\end{equation}

\begin{theorem}[Exact background-gauge negative mode]
\label{thm:s6cusp-reducible-negative-mode}
Let
\begin{align}
 \mathscr H_R^{(1)}
 &:={L^2\!\left(F_R;
 T^*F_R\otimes_{\R}\operatorname{ad}(\mathcal E_R)\right)}_{\C},
 \label{eq:s6cusp-background-hilbert-space}\\
 \mathscr D_R^{\mathrm{bg}}
 &:=\left\{a\in
 {H^2\!\left(F_R;
 T^*F_R\otimes_{\R}\operatorname{ad}(\mathcal E_R)\right)}_{\C}:
 \dd_{A_R^{\mathrm{red}}}^*a=0\right\}.
 \label{eq:s6cusp-background-sobolev-domain}
\end{align}
Relative to the unscaled \(L^2\) inner product on
\(\mathscr H_R^{(1)}\), the normalized Yang--Mills Jacobi differential
expression on \(\mathscr D_R^{\mathrm{bg}}\) is
\begin{equation}\label{eq:s6cusp-background-jacobi-operator}
 (\mathcal M_R^{\mathrm{red}}a)_\nu
 =-D^\mu D_\mu a_\nu
 -2[F^{\mathrm{red}}_{R,\nu\mu},a_\mu].
\end{equation}
For \(r=0,1\), define the smooth complex adjoint-valued one-form
\begin{equation}\label{eq:s6cusp-explicit-negative-root-mode}
 a_{r,R}^+
 :=(\pi_R^*\phi_{r,R})
 (\dd x^2-i\dd x^4)E_+.
\end{equation}
Then
\begin{align}
 \dd_{A_R^{\mathrm{red}}}^*a_{r,R}^+&=0,
 \label{eq:s6cusp-negative-mode-background-gauge}\\
 \mathcal M_R^{\mathrm{red}}a_{r,R}^+
 &=-\frac{4\pi}{\delta_R}a_{r,R}^+.
 \label{eq:s6cusp-negative-mode-eigenvalue}
\end{align}
The two \(E_+\) modes and their Hermitian-adjoint \(E_-\) partners give a real
four-dimensional negative eigenspace.  The Riesz Hessian of
\(\operatorname{YM}_R\) has eigenvalue
\begin{align}
 \lambda_R^{\mathrm{Hess}}
 &=-\frac{4\pi}{g_{\mathrm{YM}}^2\delta_R}
 \label{eq:s6cusp-negative-riesz-eigenvalue}\\
 &=\frac{4\pi}{g_{\mathrm{YM}}^2
 \operatorname{Im}\tau_R
 \left(
  \operatorname{Im}\beta_R-
  \dfrac{6(\operatorname{Im}\mu_R)^2}
  {\operatorname{Im}\tau_R}
 \right)}<0.
 \label{eq:s6cusp-negative-riesz-unchanged}
\end{align}
In particular,
\begin{equation}\label{eq:s6cusp-negative-mode-cusp}
 \lambda_R^{\mathrm{Hess}}
 =-\frac{4\pi}{g_{\mathrm{YM}}^2R^2}
 \left[
  1-\frac{2h_0-d_0}{R}+O(R^{-2})
 \right].
\end{equation}
Thus every displayed nonflat critical point is a saddle, with an explicit
negative direction whose magnitude tends to zero at the source cusp.
\end{theorem}

\begin{proof}
For completeness, expand the curvature at an arbitrary critical connection:
\[
 F_{A+ta}=F_A+t\dd_Aa+t^2a\wedge a.
\]
The second derivative of \eqref{eq:s6cusp-yang-mills-action} is
\[
 \frac1{g_{\mathrm{YM}}^2}
 \left(\lVert\dd_Aa\rVert^2
 +2\langle F_A,a\wedge a\rangle\right).
\]
Adding the background-gauge square
\(g_{\mathrm{YM}}^{-2}\lVert\dd_A^*a\rVert^2\), integrating the two
covariant derivatives by parts, and using
\([D_\mu,D_\nu]=\operatorname{ad}(F_{\mu\nu})\) gives, component by
component,
\[
 \dd_A^*\dd_A+\dd_A\dd_A^*
 =-D^\mu D_\mu-\operatorname{ad}(F_{\nu\mu}),
\]
while \(2\langle F_A,a\wedge a\rangle\) contributes the second
\(-\operatorname{ad}(F_{\nu\mu})\).  This proves the literal coefficient
\(-2\) in \eqref{eq:s6cusp-background-jacobi-operator}; it agrees with the
directly checked homogeneous-background operators
\(K_\pm=-D^2\pm2gB\) in \cite[eqs.~(1)--(2)]{Oxman2010}.

For \eqref{eq:s6cusp-explicit-negative-root-mode}, the divergence is
\[
 \dd_{A_R^{\mathrm{red}}}^*a_{r,R}^+
 =-(D_2^+-iD_4^+)\pi_R^*\phi_{r,R}=0
\]
by \eqref{eq:s6cusp-root-physical-ground-equation}.  Put
\(\mathcal B=\mathcal B_R^{\mathrm{root}}\).  Since
\(F^{\mathrm{red}}_{R,24}=(\mathcal B/2)H\), the spin-curvature term on
the \((\dd x^2,\dd x^4)E_+\) coefficients is the Hermitian matrix
\begin{equation}\label{eq:s6cusp-root-spin-matrix}
 \mathcal S_R^+
 =\mathcal B
 \begin{pmatrix}0&-2i\\2i&0\end{pmatrix}.
\end{equation}
Direct multiplication gives
\[
 \begin{pmatrix}0&-2i\\2i&0\end{pmatrix}
 \begin{pmatrix}1\\-i\end{pmatrix}
 =-2\begin{pmatrix}1\\-i\end{pmatrix}.
\]
Equation \eqref{eq:s6cusp-root-scalar-ground-energy} supplies the scalar
eigenvalue \(+\mathcal B\), whereas the displayed polarization supplies
\(-2\mathcal B\).  Their sum is
\(-\mathcal B=-4\pi/\delta_R\), proving
\eqref{eq:s6cusp-negative-mode-eigenvalue}.

Hermitian adjunction exchanges \(E_+\) with \(E_-\), while scalar complex
conjugation exchanges \(\dd x^2-i\dd x^4\) with
\(\dd x^2+i\dd x^4\) and \(\phi_{r,R}\) with
\(\overline{\phi_{r,R}}\).  Thus \((a_{r,R}^+)^\dagger\) is an
\(E_-\)-root eigenform with the same real eigenvalue.  The anti-Hermitian
combinations
\[
 a_{r,R}^+-(a_{r,R}^+)^\dagger,
 \qquad
 i\bigl(a_{r,R}^++(a_{r,R}^+)^\dagger\bigr),
 \qquad r=0,1,
\]
are four real linearly independent eigenforms.  The factor
\(g_{\mathrm{YM}}^{-2}\) in the Riesz Hessian proves
\eqref{eq:s6cusp-negative-riesz-eigenvalue};
\eqref{eq:s6cusp-unchanged-expression-delta} proves the unchanged form
\eqref{eq:s6cusp-negative-riesz-unchanged}, and
\eqref{eq:s6cusp-delta-two-term-expansion} proves the cusp expansion.
\end{proof}

The instability is therefore proved for this exact compact background; it
is not imported merely from the general Nielsen--Olesen mechanism.  Its
mathematical role here is negative: the small-action nonflat critical family
does not provide stable vacua or positive excitations above them.

\begin{proposition}[Exact relation to the native class and exceptional
kernel]
\label{prop:s6cusp-decomposable-native-relation}
In the common integral coframe,
\begin{equation}\label{eq:s6cusp-two-flux-classes}
 \xi_R=\gamma\wedge u,
 \qquad
 \eta_R=u\wedge w+6\gamma\wedge\delta.
\end{equation}
Their complete square invariants are
\begin{equation}\label{eq:s6cusp-two-flux-squares}
 \xi_R^2=0,
 \qquad
 \eta_R^2=12\gamma\wedge u\wedge w\wedge\delta.
\end{equation}
Consequently no integral lattice automorphism carries \(\eta_R\) to
\(\xi_R\).  On every exceptional elliptic kernel from
Theorem~\ref{thm:s6cusp-exceptional-elliptic-restriction}, however, the
already constructed inclusion morphism
\(\iota_{a,b,z}:E_{\tau_{\mathrm f}(z)}\hookrightarrow F_z\) gives the
exact pullbacks
\begin{equation}\label{eq:s6cusp-two-flux-kernel-pullbacks}
 \iota_{a,b,z}^*\xi_z=0,
 \qquad
 \int_{E_{\tau_{\mathrm f}(z)}}\iota_{a,b,z}^*\eta_z=6.
\end{equation}
\end{proposition}

\begin{proof}
The first formula is the definition of \(\xi_R\) and the source
identification of \((\vartheta_R^1,\ldots,\vartheta_R^4)\) with
\((\gamma,u,w,\delta)\); the second is
\eqref{eq:s6cusp-native-eta}.  Exterior multiplication gives
\eqref{eq:s6cusp-two-flux-squares}.  If an integral lattice automorphism
\(A\) obeyed \(A^*\eta_R=\xi_R\), then
\[
 0=\xi_R^2=A^*(\eta_R^2)
 =12\det(A)\gamma\wedge u\wedge w\wedge\delta,
\]
which is impossible because \(\det(A)=\pm1\).

For the kernel's positive basis
\((\widehat\delta,-\widehat\gamma_{a,b})\), direct evaluation gives
\[
 \xi_z(\widehat\delta,-\widehat\gamma
 +a\widehat u+b\widehat w)=0.
\]
This proves the first pullback.  The second is the already proved evaluation
\(\eta_z(\widehat\delta,-\widehat\gamma_{a,b})=6\) in
\eqref{eq:s6cusp-native-restriction-degree}.  Thus the two classes are
related on the same ambient object and by the same kernel morphism, with
their distinct pullbacks computed rather than inferred from their
presentations.
\end{proof}

\subsection{A subsidiary finite-cover unspooling}

A compact four-torus cannot itself be replaced by \(\R^4\): every
continuous image of \(F_R\) is compact, so no continuous map
\(F_R\to\R^4\) is surjective.  A continuous injection would be a topological
embedding because \(F_R\) is compact and \(\R^4\) is Hausdorff; invariance
of domain would make its image open, while compactness would make it
compact, which is impossible for a nonempty open subset of \(\R^4\).
Equation \eqref{eq:s6cusp-universal-cover} and the finite covers below are
the exact morphisms replacing either impossible identification.

\paragraph{Scope of the auxiliary construction.}
The native pointed limit is exactly
\eqref{eq:s6cusp-native-pointed-limit}; it is obtained from the original
marked fibres without changing their lattices.  The construction below
instead replaces \(P_R\Z^4\) by the proper sublattice
\(P_RS_M\Z^4\) and replaces the punctured base coordinate by
\(t_c=s^M\).  Its domain and codomain are therefore finite covers of the
same regular S6 fibre family.  No S12 or S24 object, and no map from such an
object into this family, occurs in the construction.  The resulting
\(\R^4\) pointed limit is a proved cover result, but it is not the separately
requested S12/S24-style globalization or doubling.

For the remainder of this subsection, \(M\ge1\) is the cover index; it is
independent of the cubical discretization index \(N\) used above.  Set
\begin{equation}\label{eq:s6cusp-cover-scaling-matrix}
 S_M:=\operatorname{diag}(1,1,M,M)
\end{equation}
in the source lattice basis
\((\widehat\gamma,\widehat u,\widehat w,\widehat\delta)\).

\begin{theorem}[Finite cover, deck group, and exact cusp monodromy]
\label{thm:s6cusp-monodromy-compatible-cover}
The sublattice
\begin{equation}\label{eq:s6cusp-cover-sublattice}
 \Lambda_M:=S_M\Z^4
 =\Z\langle
 \widehat\gamma,\widehat u,
 M\widehat w,M\widehat\delta\rangle
\end{equation}
defines the complex torus
\begin{equation}\label{eq:s6cusp-covered-fibre}
 \widetilde F_{R,M}:=
 \R^4/P_RS_M\Z^4.
\end{equation}
The identity of \(\R^4\) descends to a holomorphic covering
\begin{equation}\label{eq:s6cusp-finite-cover-map}
 c_{R,M}:\widetilde F_{R,M}\longrightarrow F_R,
 \qquad
 [x]_{P_RS_M\Z^4}\longmapsto[x]_{P_R\Z^4},
\end{equation}
with
\begin{align}
 \operatorname{Deck}(c_{R,M})
 &\cong
 \frac{P_R\Z^4}{P_RS_M\Z^4}
 \cong(\Z/M\Z)^2,
 \label{eq:s6cusp-cover-deck-group}\\
 \deg c_{R,M}&=M^2,\label{eq:s6cusp-cover-degree}\\
 \operatorname{Vol}(\widetilde F_{R,M})
 &=M^2\delta_R.
 \label{eq:s6cusp-cover-volume}
\end{align}

The clockwise cusp monodromy printed in the source is
\begin{equation}\label{eq:s6cusp-source-monodromy}
 M_0=
 \begin{pmatrix}
 1&0&0&0\\
 0&1&0&0\\
 0&1&1&0\\
 -1&0&0&1
 \end{pmatrix}
 =I+\mathcal N_0,
\end{equation}
where
\begin{equation}\label{eq:s6cusp-source-nilpotent}
 \mathcal N_0(a,b,c,d)^{\mathsf T}
 =(0,0,b,-a)^{\mathsf T},
 \qquad
 \mathcal N_0^2=0.
\end{equation}
After the cyclic base change
\begin{equation}\label{eq:s6cusp-cover-base-change}
 r_M:\Delta_s^*\longrightarrow\Delta_{t_c}^*,
 \qquad t_c=s^M,
\end{equation}
the pulled-back monodromy preserves \(\Lambda_M\) and, in its displayed
basis, is again \(M_0\):
\begin{align}
 M_0^M&=I+M\mathcal N_0,
 \label{eq:s6cusp-monodromy-power}\\
 M_0^M\Lambda_M&=\Lambda_M,
 \label{eq:s6cusp-monodromy-preserves-sublattice}\\
 S_M^{-1}M_0^MS_M&=M_0.
 \label{eq:s6cusp-monodromy-conjugacy}
\end{align}
Consequently the fibres \(\widetilde F_{R,M}\) assemble into a holomorphic
torus family over \(\Delta_s^*\), and the maps \(c_{R,M}\) assemble into a
fibrewise \(M^2\)-sheeted holomorphic covering of the pulled-back source
family.
\end{theorem}

\begin{proof}
The inclusion \(S_M\Z^4\subset\Z^4\) and invertibility of \(P_R\) prove
that \eqref{eq:s6cusp-finite-cover-map} is well defined.  Its fibre over a
point is the quotient of \(P_R\Z^4\) by \(P_RS_M\Z^4\).  Multiplication by
\(P_R^{-1}\) identifies this quotient with
\(\Z^4/S_M\Z^4\), and the explicit homomorphism
\[
 [(n_1,n_2,n_3,n_4)]
 \longmapsto([n_3],[n_4])
\]
has kernel \(S_M\Z^4\) and image \((\Z/M\Z)^2\).  This proves
\eqref{eq:s6cusp-cover-deck-group} and
\eqref{eq:s6cusp-cover-degree}.  The covolume is
\[
 |\det(P_RS_M)|=|\det P_R|\det S_M=M^2\delta_R,
\]
proving \eqref{eq:s6cusp-cover-volume}.  Since the period inclusions are
complex translations in \(\C^2\), the covering is holomorphic.

Equations \eqref{eq:s6cusp-source-monodromy} and
\eqref{eq:s6cusp-source-nilpotent} are the coordinate form of
Definition~2.3 and Lemmas~2.4 and~2.6 of the exact S6 transcription
\cite{S6Internal}.  The binomial formula terminates because
\(\mathcal N_0^2=0\), giving
\eqref{eq:s6cusp-monodromy-power}.  On the four generators of
\(\Lambda_M\),
\begin{align*}
 M_0^M\widehat\gamma
 &=\widehat\gamma-M\widehat\delta,&
 M_0^M\widehat u
 &=\widehat u+M\widehat w,\\
 M_0^M(M\widehat w)&=M\widehat w,&
 M_0^M(M\widehat\delta)&=M\widehat\delta.
\end{align*}
These equations prove invariance.  Applying \(S_M^{-1}\) to the four
displayed images gives exactly the four columns of \(M_0\), proving
\eqref{eq:s6cusp-monodromy-conjugacy}.

One positive loop around \(s=0\) maps to \(M\) loops around \(t_c=0\),
so the pulled-back lattice local system has monodromy \(M_0^M\).
The proved invariant sublattice therefore defines a sub-local-system.
Quotienting the pulled-back \(\C^2\)-bundle by its translations constructs
the asserted holomorphic torus family, and inclusion into the full
pulled-back lattice gives the family map whose fibre is
\eqref{eq:s6cusp-finite-cover-map}.
\end{proof}

For \(M=2\), the cover in
\eqref{eq:s6cusp-finite-cover-map} has four sheets.  Together with the
twofold base change \(t_c=s^2\), equation
\eqref{eq:s6cusp-monodromy-conjugacy} is the exact four-sheet realization
of the proposed two-cycle doubling; no tetrahedral or Fable identification
is used in this lattice statement.

\begin{theorem}[Systole bound and pointed Euclidean four-space limit]
\label{thm:s6cusp-covered-R4-limit}
Every period of \(\widetilde F_{R,M}\), in the reordered covering
coordinates \((x^1,x^3;x^2,x^4)\), has the unique form
\begin{equation}\label{eq:s6cusp-covered-period-block-form}
 \lambda_{R,M}(n,m)=
 \begin{pmatrix}
 A_Rn+Mm\\
 B_Rn
 \end{pmatrix},
 \qquad n,m\in\Z^2,
\end{equation}
with \(A_R,B_R\) as in
\eqref{eq:s6cusp-AR-matrix} and
\eqref{eq:s6cusp-BR-matrix}.  Its systole and injectivity radius obey
\begin{align}
 \operatorname{sys}(\widetilde F_{R,M})
 &\ge\min\{R-C_*,M\},
 \label{eq:s6cusp-covered-systole-bound}\\
 \operatorname{inj}(\widetilde F_{R,M})
 &\ge\frac12\min\{R-C_*,M\}.
 \label{eq:s6cusp-covered-injectivity-bound}
\end{align}
For every pair of sequences \(R_j\to\infty\) and \(M_j\to\infty\),
\begin{equation}\label{eq:s6cusp-covered-pointed-limit}
 (\widetilde F_{R_j,M_j},[0])
 \longrightarrow(\R^4,0)
\end{equation}
in pointed smooth, pointed local-isometric, and pointed
Gromov--Hausdorff senses.

Along the radial path \(t_c=e^{-2\pi R}\), the positive real lift under
\eqref{eq:s6cusp-cover-base-change} is
\begin{equation}\label{eq:s6cusp-covered-base-radial-coordinate}
 s=e^{-2\pi R/M}.
\end{equation}
Hence the points also approach the cusp \(s=0\) of the varying covered
bases exactly when \(R_j/M_j\to\infty\).  The explicit sequence
\[
 R_j=j^2,\qquad M_j=j
\]
satisfies all three divergence conditions and obeys
\[
 \operatorname{sys}(\widetilde F_{j^2,j})
 \ge j
\]
for all sufficiently large \(j\).
\end{theorem}

\begin{proof}
The first two columns of \(P_RS_M\) are unchanged and the last two are
multiplied by \(M\).  Separating the top coordinates \(x^1,x^3\) from the
bottom coordinates \(x^2,x^4\) gives
\eqref{eq:s6cusp-covered-period-block-form}; uniqueness follows from
invertibility of \(P_RS_M\).

For a nonzero period with \(n\ne0\),
\eqref{eq:s6cusp-BR-lower-bound} gives
\[
 \lVert\lambda_{R,M}(n,m)\rVert
 \ge\lVert B_Rn\rVert
 \ge R-C_*.
\]
For \(n=0\), nonzero means \(m\ne0\), and
\[
 \lVert\lambda_{R,M}(0,m)\rVert=M\lVert m\rVert\ge M.
\]
Taking the minimum proves \eqref{eq:s6cusp-covered-systole-bound}; the
flat-torus identity
\(\operatorname{inj}=\operatorname{sys}/2\) proves
\eqref{eq:s6cusp-covered-injectivity-bound}.

Fix \(L>0\).  For all large \(j\), the systole exceeds \(4L\).  Let
\[
 \widetilde\upsilon_{R_j,M_j}:
 \R^4\longrightarrow\widetilde F_{R_j,M_j}
\]
be the quotient map.  If \(x,y\in B_L(0)\) and \(\lambda\) is a nonzero
period, then
\[
 \lVert x-y+\lambda\rVert
 \ge\lVert\lambda\rVert-\lVert x-y\rVert
 >4L-2L\ge\lVert x-y\rVert.
\]
Thus the zero period uniquely realizes the quotient distance between these
points, and the restriction of \(\widetilde\upsilon_{R_j,M_j}\) to
\(B_L(0)\) is an isometric embedding.  Conversely, every quotient point
at distance less than \(L\) from \([0]\) has a distance-minimizing lift in
\(B_L(0)\), because the period lattice is discrete.  The radius-\(L\)
pointed balls are therefore isometric for all large \(j\).  Since \(L\)
was arbitrary and the covering maps are local Euclidean isometries and
holomorphic maps from \(\C^2\), all three modes of convergence in
\eqref{eq:s6cusp-covered-pointed-limit} follow.

Equation \eqref{eq:s6cusp-covered-base-radial-coordinate} is the unique
positive real solution of \(s^M=e^{-2\pi R}\).  It tends to zero exactly
when \(R/M\to\infty\).  Substitution of \(R_j=j^2,M_j=j\) into
\eqref{eq:s6cusp-covered-systole-bound} finishes the proof.
\end{proof}

\begin{theorem}[Exact field pullback and added momentum sectors]
\label{thm:s6cusp-covered-dual-modes}
The scalar characters of the covered fibre are
\begin{equation}\label{eq:s6cusp-covered-characters}
 \widetilde\chi_{R,M,k}([x])
 :=
 \exp\left(
 2\pi i\left\langle
 P_R^{-\mathsf T}S_M^{-1}k,x
 \right\rangle
 \right),
 \qquad k\in\Z^4,
\end{equation}
and the two mutually inverse dual-lattice maps are
\[
 k\longmapsto P_R^{-\mathsf T}S_M^{-1}k,
 \qquad
 q\longmapsto S_MP_R^{\mathsf T}q.
\]
Pullback under the finite cover is exactly
\begin{equation}\label{eq:s6cusp-covered-character-pullback}
 c_{R,M}^*\chi_{R,\ell}
 =\widetilde\chi_{R,M,S_M\ell}.
\end{equation}
Consequently its image consists precisely of the characters whose labels
satisfy
\begin{equation}\label{eq:s6cusp-old-mode-congruences}
 k_3\equiv0\pmod M,\qquad k_4\equiv0\pmod M.
\end{equation}
The scalar spectrum is
\begin{equation}\label{eq:s6cusp-covered-scalar-spectrum}
 \operatorname{Spec}(\widetilde\Delta_{R,M})
 =
 \left\{
 4\pi^2
 \lVert P_R^{-\mathsf T}S_M^{-1}k\rVert^2:
 k\in\Z^4
 \right\}.
\end{equation}
After harmonic forms are removed, the normalized flat Yang--Mills Hessian
has the same positive eigenvalue set,
\begin{equation}\label{eq:s6cusp-covered-gauge-spectrum}
 \operatorname{Spec}^+(\widetilde{\mathcal L}_{R,M})
 =
 \left\{
 4\pi^2
 \lVert P_R^{-\mathsf T}S_M^{-1}k\rVert^2:
 k\in\Z^4\setminus\{0\}
 \right\},
\end{equation}
with three complex transverse polarizations per \(k\) and per
\(\mathfrak g_{\C}\)-direction.  Along every sequence in
\eqref{eq:s6cusp-covered-pointed-limit},
\begin{equation}\label{eq:s6cusp-covered-hessian-gap-limit}
 \inf\operatorname{Spec}^+(\widetilde{\mathcal L}_{R_j,M_j})
 \longrightarrow0.
\end{equation}
\end{theorem}

\begin{proof}
A covector pairs integrally with \(P_RS_M\Z^4\) exactly when
\(S_MP_R^{\mathsf T}q\in\Z^4\), which proves the displayed inverse maps
and the well-definedness of \eqref{eq:s6cusp-covered-characters}.  Both
sides of \eqref{eq:s6cusp-covered-character-pullback}, evaluated on a
representative \(x\), equal
\[
 \exp\left(2\pi i
 \langle P_R^{-\mathsf T}\ell,x\rangle\right).
\]
The image labels are \(S_M\ell=(\ell_1,\ell_2,M\ell_3,M\ell_4)\), which
proves both directions of \eqref{eq:s6cusp-old-mode-congruences}.
Direct differentiation proves
\eqref{eq:s6cusp-covered-scalar-spectrum}; the quotient calculation in
Theorem~\ref{thm:s6cusp-flat-gauge-hessian}, with \(P_R\) replaced by
\(P_RS_M\), proves \eqref{eq:s6cusp-covered-gauge-spectrum} and its
multiplicity.

For an explicit upper bound, the labels \(e_1\) and \(e_3\) give
\begin{equation}\label{eq:s6cusp-covered-gap-upper-bound}
 \inf\operatorname{Spec}^+(\widetilde{\mathcal L}_{R,M})
 \le4\pi^2\min\left\{
 \lVert P_R^{-\mathsf T}e_1\rVert^2,
 \frac1{M^2}\lVert P_R^{-\mathsf T}e_3\rVert^2
 \right\}.
\end{equation}
The first term tends to zero by
\eqref{eq:s6cusp-e1-exact}.  For the second, the full dual formula gives
\begin{equation}\label{eq:s6cusp-third-dual-coordinate}
 P_R^{-\mathsf T}e_3=
 \left(
 1,\frac{p_Ru_R-6a_Rm_R}{\delta_R},0,
 \frac{6a_RT_R-6m_Ru_R}{\delta_R}
 \right),
\end{equation}
which remains bounded because \(a_R,u_R,m_R\) are bounded,
\(p_R,T_R=O(R)\), and \(\delta_R=R^2+O(R)\).
Thus either entry on the right of
\eqref{eq:s6cusp-covered-gap-upper-bound} proves the limit, and together
they exhibit the momentum scales supplied by the \(R\)- and \(M\)-directions.
\end{proof}

The modes violating \eqref{eq:s6cusp-old-mode-congruences} are not pullbacks
of fields on \(F_R\).  Their third and fourth label components enter
\(S_M^{-1}k\) as \(k_3/M,k_4/M\), so the finite cover supplies genuinely
new fractional momenta in the two directions that remain compact in the
native limit.  This is the exact field-space map behind the geometric
unspooling.

\begin{proposition}[The alternative doubled quotient and its exact kernel]
\label{prop:s6cusp-doubled-quotient}
Let
\[
 D_R:=\Sigma_R\times\Sigma_R
 =\R^4/
 \begin{pmatrix}B_R&0\\0&B_R\end{pmatrix}\Z^4.
\]
The map
\begin{equation}\label{eq:s6cusp-doubled-quotient-map}
 \Pi_R:F_R\times F_R\longrightarrow D_R,
 \qquad
 (x,y)\longmapsto(\pi_R(x),\pi_R(y))
\end{equation}
is a surjective homomorphism with
\begin{equation}\label{eq:s6cusp-doubled-quotient-kernel}
 \ker\Pi_R=K_R\times K_R.
\end{equation}
Moreover,
\[
 \operatorname{sys}(D_R)\ge R-C_*,
 \qquad
 (D_R,[0])\longrightarrow(\R^4,0)
\]
in the same three pointed senses as
\eqref{eq:s6cusp-covered-pointed-limit}.
\end{proposition}

\begin{proof}
Surjectivity and the kernel formula are the product of the proved exact
sequence \eqref{eq:s6cusp-native-exact-sequence} with itself.  Every
nonzero period of \(\Sigma_R\) is \(B_Rn\) for \(n\ne0\), so
\eqref{eq:s6cusp-BR-lower-bound} gives
\(\operatorname{sys}(\Sigma_R)\ge R-C_*\).  The product lattice has the
same shortest length.  The expanding-systole ball argument in the proof of
Theorem~\ref{thm:s6cusp-covered-R4-limit} proves the pointed limit.
\end{proof}

The morphism \eqref{eq:s6cusp-doubled-quotient-map} removes the four-real-
dimensional compact subgroup in
\eqref{eq:s6cusp-doubled-quotient-kernel}; the finite-cover construction
instead retains those directions and increases their periods.  These two
presentations are therefore related by the exact quotient map above, not
by an identification.

\begin{proposition}[The local Wick map after four-direction unspooling]
\label{prop:s6cusp-local-wick-map}
On the complexification of the Euclidean pointed limit define
\[
 W:\R^{1,3}\longrightarrow\C^4,\qquad
 W(t,x_1,x_2,x_3)=(it,x_1,x_2,x_3).
\]
Then
\begin{equation}\label{eq:s6cusp-wick-pullback}
 W^*\left(\dd z_0^2+\dd z_1^2+\dd z_2^2+\dd z_3^2\right)
 =-\dd t^2+\dd x_1^2+\dd x_2^2+\dd x_3^2.
\end{equation}
The map is injective, with image
\(\{(z_0,z_1,z_2,z_3):z_0\in i\R,\ z_1,z_2,z_3\in\R\}\), and inverse
\((z_0,z_1,z_2,z_3)\mapsto(-iz_0,z_1,z_2,z_3)\) on that image.
\end{proposition}

\begin{proof}
One has \(W^*\dd z_0=i\,\dd t\) and
\(W^*\dd z_j=\dd x_j\) for \(j=1,2,3\).  Squaring these four retained
terms gives \eqref{eq:s6cusp-wick-pullback}; the displayed image and inverse
prove the remaining assertions.  The calculation changes the real slice
of the complexified bilinear form.  It does not remove a lattice period,
which is why it is placed after the finite-cover limit rather than used as
a substitute for it.
\end{proof}
